\chapter{Major open problems}
\label{chap:major-open-problems}
\index{open problem!major list}
\index{research frontier}

This chapter gives a source-dated list of the main problems left by the
theories developed in this book.  Its cutoff is 24 August 2026.  The list is
selective rather than encyclopedic: a problem is included when solving it
would remove a structural boundary between two parts of the subject, extend a
core existence theorem to its natural domain, or settle a repeatedly used
regularization or compactness gap.

Every problem is accompanied by three safeguards.  First, its domain is part
of its statement.  Second, the strongest known positive cases and negative
obstructions are recorded before any proposed route.  Third, each proposed
solution is submitted to the verification protocol of
\Cref{chap:verified-extensions}.  The status words used below are:

\begin{description}[style=nextline,leftmargin=1.5em,labelindent=0pt,
  itemsep=.45\baselineskip,font=\normalfont\bfseries]
\item[Verified] the displayed statement has a complete proof or a controlling
primary-source theorem;
\item[Disproved] the displayed signature is impossible, with an explicit
counterexample or formal obstruction;
\item[Conditional] the conclusion follows from a precisely stated missing
input;
\item[Open] the attempted proof and counterexample search did not close the
displayed signature.
\end{description}

The literal inventory is exact.  This chapter contains \(21\) problem
environments.  Together with the separate projective-rational full-cone
problem \Cref{prob:pseff-projective-rational-full-link} in the preceding
chapter, the book-wide frontier inventory contains \(22\) problem
environments.  Each of those \(22\) has a genuinely open residual domain as
stated; no environment is wholly resolved, wholly disproved, or merely
conditional.  Thus the whole-environment count is
\[
 \text{resolved }0,\qquad
 \text{disproved }0,\qquad
 \text{conditional }0,\qquad
 \text{genuinely open }22.
\]
Verified, disproved, and conditional assertions below are boundary cases or
subclaims inside those open environments.  A source question that has been
resolved or disproved in its entirety is recorded for chronology but is not
silently counted as an active problem.

An open verdict is not a placeholder for a proof.  The accompanying attempted
route is included so that the next reader begins at the actual obstruction
rather than repeating a known invalid argument.  New verified results and
counterexamples are isolated in complete proof blocks either here or in
\Cref{chap:verified-extensions}; the complete machine-verification receipts
are recorded in the editorial ledger.

\section{Local positivity and the CLD--Gubler--Rabinoff frontier}

\subsection{Pulling positive PL currents back to curves}

\begin{problem}\label{prob:frontier-current-pullback}
Let \(K\) be complete and nontrivially valued, \(X\) a good separated strict
\(K\)-analytic space, and \(h\) a real-valued real PL function.  Does
\[
                         \dprimeddprime[h]\geq0
\]
on the relative interior imply that \(h|_{\operatorname{Int}(X)}\) is
residue-nef, equivalently model semipositive?  Equivalently, one may assume
from the outset that \(X\) is boundaryless and ask for semipositivity on
\(X\).
\end{problem}

The converse direction is known: semipositive real PL functions have positive
current on the relative interior by \cite[Theorem~7.14]{GJR26}.  Model
semipositivity is equivalent to the functorial classical curve test by
\cite[Theorem~4.3]{GR25Classical}, and it is detected on sufficiently small
analytic curves by \cite[Theorem~5.3]{GR25PL}.  Thus the missing step is not a
one-dimensional positivity theorem; it is a positivity-preserving restriction
or slicing operation for the singular current \([h]\) along every curve in
the criterion.  This obstruction is explicitly recorded in
\cite[Remark~1.8]{GR25PL}.

The completed Rethlas audit gives four separate verdicts.  The main analytic
converse is still \emph{open}.  It is deliberately not folded into the
verified equivalence \eqref{eq:verified-pl-psh-equivalence}.  The exact
conditional reduction is the following restriction statement:
\[
 \begin{gathered}
 \dprimeddprime[h]\geq0
 \quad\Longrightarrow\quad
 \dprimeddprime[h\circ f]\geq0,\\
 f:C\longrightarrow\operatorname{Int}(X)
 \quad\text{any sufficiently small criterion-curve map}.
 \end{gathered}
\tag{R}\label{eq:frontier-restriction-r}
\]
where curve test functions are compactly supported away from the boundary.
Assuming \eqref{eq:frontier-restriction-r}, positivity on each curve says
that every outgoing-slope sum is nonnegative.  The local curve-detection
theorem \cite[Theorem~5.3]{GR25PL}, followed by
\cite[Theorem~4.3]{GR25Classical}, then gives model semipositivity.  Together
with \cite[Theorem~7.14]{GJR26}, this proves the desired equivalence
conditionally on \eqref{eq:frontier-restriction-r}.

The naive tropical substitute for \eqref{eq:frontier-restriction-r} is
\emph{disproved}.  Let
\[
 \Sigma=P_{xy}\cup P_{xz}\subset\mathbf R^3,
 \qquad P_{xy}=\{z=0\},\qquad P_{xz}=\{y=0\},
\]
give both planes weight one, and let \(L\) be their common \(x\)-axis.  For
\(a>b>0\), the continuous PL function
\[
 \varphi|_{P_{xy}}=a|y|,
 \qquad
 \varphi|_{P_{xz}}=-b|z|
\]
has aggregate divisor
\[
             \operatorname{div}_{\Sigma}(\varphi)=2(a-b)[L]\geq0,
\]
but its restriction to the balanced \(z\)-axis \(C=\{x=y=0\}\) has
\[
             \operatorname{div}_{C}(\varphi|_C)=-2b[0]<0.
\]
Thus positivity after coarse branch aggregation does not imply positivity on
a selected tropical curve.  In the evident analytic realization by
\(\{z=1\}\) and \(\{y=1\}\) inside a torus, however, a refined test on the
second branch already detects the negative jump.  The ambient CLD current is
therefore not positive.  This is an obstruction to a proof strategy, not a
Berkovich counterexample to the problem.

A precise sufficient input for \eqref{eq:frontier-restriction-r} would be a
positive Thom-form approximation.  Put \(U=\operatorname{Int}(X)\).  For
every criterion curve \(f:C\to U\)
and every nonnegative compactly supported smooth curve test function
\(\psi\), require strongly positive compactly supported ambient
\((n-1,n-1)\)-forms \(\alpha_j\) such that, with analytic multiplicities,
\[
 \bigl\langle\dprimeddprime[h],\alpha_j\bigr\rangle
 \longrightarrow
 \bigl\langle\dprimeddprime[h\circ f],\psi\bigr\rangle.
\]
Ambient positivity would make the limit nonnegative.  No such approximation
theorem is presently available in the required singular and excess
intersection generality.  By contrast, if a finite ambient PL potential is
convex, its restriction to any balanced tropical curve is subharmonic:
sublinearity of its directional derivative and
\(\sum_i m_i u_i=0\) give
\(\sum_i m_iD\varphi(u_i)\geq D\varphi(0)=0\).  This verifies the convex
boundary case but does not settle the singular-current problem.

\subsection{Regularizing arbitrary classical psh functions}

\begin{problem}\label{prob:frontier-classical-regularization}
On which higher-dimensional good strict analytic spaces is every classically
psh function locally a decreasing limit of semipositive real PL, locally
Fubini--Study, weakly smooth psh, or CLD-smooth psh functions?
\end{problem}

The four targets must be separated.  Real PL functions on the relative
interior satisfy the three-way comparison of
\Cref{thm:verified-pl-psh-equivalence}.  On smooth boundaryless curves,
\cite[Proposition~15.11]{GJR26} gives decreasing weakly smooth
regularization locally on each relatively compact open subset.  Wanner's
CLD-smooth regularization has the additional
residue-field or curve hypotheses stated in \cite[Corollary~5.4]{Wanner19};
weakly smooth harmonic functions need not be CLD-smooth.

The attempted higher-dimensional proof combines local formal models,
regularized maxima, and harmonic tropicalizations.  Curvewise approximation
does not by itself glue: the approximants must be compatible on overlapping
tropical charts and after scalar extension.  Nor may a nonmetrizable net be
replaced by a sequence without a countable extraction argument.  The general
problem is \emph{open}; the curve and PL cases are verified boundary cases.

\subsection{Upper-semicontinuous regularization of local envelopes}

\begin{problem}\label{prob:frontier-local-envelope}
For an upper-semicontinuous obstacle
\(f:U\to\mathbf R\cup\{-\infty\}\) on a good strict analytic domain \(U\),
let
\[
 P_U^{\mathrm{raw}}f
 :=\sup\{g:g\text{ is local FS or semipositive real PL on }U,\ g\leq f\}.
\]
When is \((P_U^{\mathrm{raw}}f)^*\) classically psh, CLD-v2 psh, or
continuous?
\end{problem}

There is a sharp formal reduction.  If the raw supremum is already usc, then
it is classically psh: direct finite subsets of the family by inclusion, take
their maxima, and apply the upper-semicontinuous net-closure theorem
\cite[Theorem~6.5]{GR25Classical}.  The missing assertion is that usc
regularization commutes with restriction to every analytic curve after every
scalar extension.  Assuming that commutation, together with the corresponding
upper-envelope theorem for subharmonic functions on curves, proves the
classical psh conclusion.

The Rethlas search begins with curves, polystable charts, and domains carrying
a relatively ample formal model, before testing relative-boundary points.
Continuity needs barriers or a local capacity estimate in addition to psh
closure.  The raw-usc case is \emph{verified}; the unrestricted regularized
envelope remains \emph{open}.

\subsection{A common singular local Monge--Ampère product}

\begin{problem}\label{prob:frontier-local-ma}
Construct a local product that simultaneously extends the continuous CLD
Bedford--Taylor product, the bounded formal BFJ measure on a normalized
valuation link, and the largest justified class of classically psh functions.
It should have a stated boundary datum, model independence, monotone
continuity, locality, and functoriality.
\end{problem}

The two existing endpoints have different signatures.  CLD works on analytic
opens with tropical charts and compactly supported test forms; the local BFJ
operator of \Cref{thm:local-formal-ma} takes \(n-1\) potentials on the link
\(\nu(\mathfrak m)=1\) and uses exceptional intersection theory.  Classical
psh positivity, by itself, constructs no wedge product.  The following
result identifies the two products on one explicit model locus and then on a
jointly realizable bounded monotone closure.  It does not identify either
product for an arbitrary curve-test psh function.

% RETHLAS-VERIFIED: local-singular-ma-comparison; receipt 20260814T170257Z_923e70de7a01
Let \((X,0)\) be a smooth complex analytic germ of dimension \(n\geq2\), put
\(R=\mathcal O_{X,0}\), let \(\mathfrak m\) be its maximal ideal, and set
\(d=n-1\).  Its normalized valuation link is
\begin{equation}
 V_{\mathfrak m}
 :=\{v:R\longrightarrow[0,+\infty]\mid
       v\text{ is centered at }0,\ v(\mathfrak m)=1\}.
 \label{eq:frontier-local-ma-link}
\end{equation}
For a prime divisor \(E\) above \(0\), write
\begin{equation}
 b_E:=\ord_E(\mathfrak m),
 \qquad v_E:=b_E^{-1}\ord_E.
 \label{eq:frontier-local-ma-divisorial-normalization}
\end{equation}
Fix \(K=\mathbf C((t))\), normalized by \(|t|=\rme^{-1}\), and regular
parameters \(z_1,\ldots,z_n\) of \(R\).  We choose a closed good, possibly
non-strict, \(K\)-analytic polydisc \(B\) of polyradius \(\rho\), where
\(|t|<\rho<1\).  Taylor series of the chosen complex analytic germs converge
on \(B\), and the shell below has a collar in \(B\).  For an
\(\mathfrak m\)-primary ideal
\(\mathfrak a=(f_1,\ldots,f_N)\), put
\begin{equation}
 u_{\mathfrak a}:=\log\max_j|f_j|,
 \qquad
 u_{\mathfrak m}:=\log\max_i|z_i|,
 \qquad
 \chi_{\mathfrak m}:=\max\{u_{\mathfrak m},-1\}.
 \label{eq:frontier-local-ma-model-functions}
\end{equation}
Changing generators leaves the associated model metric unchanged after a
common admissible refinement.  Set
\begin{equation}
S_{\mathfrak m}:=\{u_{\mathfrak m}=-1\}.
\label{eq:frontier-local-ma-shell}
\end{equation}
In this subsection an \emph{algebraizable log model} means the
analytification, after shrinking the germ, of a projective modification of a
Zariski neighborhood \(U\subset\mathbf A^n_{\mathbf C}\), with the relevant
divisors and line bundles descending to that algebraic model.

For model psh potentials \(u_1,\ldots,u_d\), let
\(T=\dprimeddprime u_1\wedge\cdots\wedge\dprimeddprime u_d\).  We define
their swept boundary product by
\begin{equation}
 \begin{split}
 \MA^{\partial}_{\mathrm{CLD},\mathfrak m}
       (u_1,\ldots,u_d)
 :={}&T\wedge\dprimeddprime\chi_{\mathfrak m}\\
 &-\one_{\{u_{\mathfrak m}>-1\}}
       \bigl(T\wedge\dprimeddprime u_{\mathfrak m}\bigr),
 \end{split}
 \label{eq:frontier-local-ma-swept-product}
\end{equation}
where the restricted measure in the second line is extended by zero.  The
two products are the continuous Bedford--Taylor products in the revised CLD
psh sheaf.  The subtraction and the level \(-1\) are part of the boundary
datum; in particular, \eqref{eq:frontier-local-ma-swept-product} is neither a
probability normalization nor the raw \(d\)-fold product of restrictions to
the shell.

The definition is meaningful for the model potentials above.  Indeed, if
\(\mathfrak a\) is \(\mathfrak m\)-primary, then
\(\mathfrak m^p\subset\mathfrak a\subset\mathfrak m^q\) for suitable
positive integers \(p,q\), so \(u_{\mathfrak a}\) is bounded on a collar of
\(S_{\mathfrak m}\).  On \(\{u_{\mathfrak m}>-1\}\), locality gives
\(\dprimeddprime\chi_{\mathfrak m}
=\dprimeddprime u_{\mathfrak m}\), and the two terms cancel.  On
\(\{u_{\mathfrak m}<-1\}\), the cutoff is constant and both terms vanish.
Thus \eqref{eq:frontier-local-ma-swept-product} is supported on the shell.
This argument does not assert that
\(\dprimeddprime\chi_{\mathfrak m}\) itself is shell-supported; that
assertion is false in dimension greater than one.  The reference factor is
essential: \(u_{\mathfrak m}|_{V_{\mathfrak m}}=-1\), whose intrinsic
\(\dprimeddprime\) is zero, whereas the local BFJ measure of this datum has
mass one.

\begin{lemma}\label{lem:frontier-local-ma-gauss}
There are maps
\[
 \iota:V_{\mathfrak m}\hookrightarrow S_{\mathfrak m},
 \qquad
 p:S_{\mathfrak m}\longrightarrow V_{\mathfrak m},
 \qquad p\circ\iota=\mathrm{id},
\]
such that, for every \(\mathfrak m\)-primary ideal \(\mathfrak a\),
\begin{equation}
             u_{\mathfrak a}(\iota(v))=-v(\mathfrak a).
 \label{eq:frontier-local-ma-gauss-ideal}
\end{equation}
If \(\pi:Y\to X\) is an algebraizable log resolution of finitely many
\(\mathfrak m\)-primary ideals and
\(\mathfrak m\mathcal O_Y=\mathcal O_Y(-D)\), then every face
\begin{equation}
 \Delta(Y,D)=
 \left\{\sum_{j=1}^r w_jE_j:\ w_j\geq0,
                       \ \sum_{j=1}^r b_jw_j=1\right\}
 \label{eq:frontier-local-ma-dual-face}
\end{equation}
of its normalized dual complex is carried by \(\iota\) onto a polyhedral
skeleton in a formal model of the shell.
\end{lemma}

\begin{proof}
For \(v\in V_{\mathfrak m}\), define its Gauss extension on convergent
Laurent series by
\begin{equation}
 \widetilde v\!\left(\sum_{j\geq j_0}t^jf_j\right)
       :=\min_j\{j+v(f_j)\}.
 \label{eq:frontier-local-ma-gauss-extension}
\end{equation}
The minimum is attained because \(v(f_j)\geq0\) and \(j\to+\infty\).
Multiplicativity follows by passing to the associated graded domain:
the initial form of the Laurent polynomial is a nonzero Laurent polynomial
in a new indeterminate of degree one over \(\operatorname{gr}_vR\), so the
initial forms of two nonzero factors cannot multiply to zero.  Moreover,
\(v(z_i)\geq1\), hence the resulting seminorm has
\(|z_i|\leq|t|<\rho\).  It is bounded on the affinoid algebra of \(B\) and
therefore extends by continuity from the dense Laurent-polynomial
subalgebra.
Its restriction to \(R\) is \(v\), and
\(\min_i\widetilde v(z_i)=v(\mathfrak m)=1=\widetilde v(t)\), so it defines
a point \(\iota(v)\) of the shell.  Conversely, restriction of the
semivaluation of a shell point to \(R\) gives an element of
\(V_{\mathfrak m}\); this is \(p\).  Formula
\eqref{eq:frontier-local-ma-gauss-extension} gives
\(p\circ\iota=\mathrm{id}\), and
\eqref{eq:frontier-local-ma-gauss-ideal} follows by taking the minimum over
the generators of \(\mathfrak a\).

On a stratum \(E_1\cap\cdots\cap E_r\), choose local equations
\(x_j=0\).  A monomial valuation has \(v(x_j)=w_j\).  Since a local
generator of \(\mathfrak m\mathcal O_Y\) has divisor
\(D=\sum b_jE_j\), normalization is exactly
\(\sum b_jw_j=1\).

For the promised formal realization, take the formal completion of
\(Y\times\Spec K^\circ\) over the chosen neighborhood and the normalized
admissible blowup of
\[
              (t,\mathfrak m\mathcal O_Y)
              =(t,\mathcal O_Y(-D)).
\]
Locally this is the two-generated ideal \((t,s_D)\), and it does not alter
the generic fibre of the already chosen modification because \(t\) is
invertible there.  Near the stratum,
\(s_D=\varepsilon x_1^{b_1}\cdots x_r^{b_r}\) with \(\varepsilon\) a unit,
so the exceptional skeleton at \(v(t)=1\) is precisely
\eqref{eq:frontier-local-ma-dual-face}.  Toroidal blowups of strata produce
common subdivisions.  Formula \eqref{eq:frontier-local-ma-gauss-extension}
identifies their skeleton points with the Gauss extensions, and every
principalized ideal function is integral affine on each face.
Such an algebraizable simultaneous resolution always exists for the finite
ideal data used here, by the polynomialization argument in
\Cref{lem:frontier-local-ma-common-model-core}.
\end{proof}

\begin{lemma}\label{lem:frontier-local-ma-common-model-core}
Let \(\mathfrak a\) be \(\mathfrak m\)-primary and suppose that
\(\mathfrak a\mathcal O_Y=\mathcal O_Y(-A)\) on a principalization.  Then
\(u_{\mathfrak a}\) is a continuous semipositive model function, hence a
member of the revised CLD psh sheaf on a collar of the shell; its pullback to
every analytic curve is subharmonic; and its restriction to the Gauss
section is the BFJ model formal psh function
\(g_{\mathfrak a}(v)=-v(\mathfrak a)\).  The same conclusion holds for a
rational semipositive model function associated with a relatively nef
exceptional \(\mathbf Q\)-Cartier b-divisor carried by an algebraizable
common log model.
\end{lemma}

\begin{proof}
We first record why an admissible algebraic formal model is available for
the ideal locus, even though the original germs are analytic.  Identify
\(R\) with \(\mathbf C\{z_1,\ldots,z_n\}\).  Choose \(N\) such that
\(\mathfrak m^N\subset\mathfrak m\mathfrak a\), and replace each generator
by its Taylor polynomial modulo \(\mathfrak m^N\).  If \(\mathfrak b\) is
the ideal generated by these polynomials, then
\[
 \mathfrak b\subset\mathfrak a
 \quad\text{and}\quad
 \mathfrak a\subset\mathfrak b+\mathfrak m^N
                 \subset\mathfrak b+\mathfrak m\mathfrak a.
\]
Nakayama's lemma applied to \(\mathfrak a/\mathfrak b\) gives
\(\mathfrak a=\mathfrak b\).  A finite collection of primary ideals can be
treated with one \(N\).  This replacement also preserves the analytic
weight.  Indeed, the transition coefficients between two finite generating
sets lie in \(R\), hence have norm at most one on the polydisc of radius
\(\rho<1\).  Applying the ultrametric inequality in both directions gives
equality of the two maxima of generator norms.  We may therefore choose an
algebraic simultaneous log resolution over a Zariski neighborhood of the
origin, isomorphic away from the origin, and later take its \(t\)-adic formal
completion.

On that resolution write locally \(f_j=s_Ah_j\).  The \(h_j\) generate the
unit ideal.  On every formal affine chart, all \(h_j\) and the coefficients
of a relation \(1=\sum_jc_jh_j\) are power bounded.  The ultrametric
inequality consequently gives
\[
                    \max_j|h_j|=1,
 \qquad
                    \log\max_j|f_j|=\log|s_A|.
\]
Thus the weight is exactly the model weight of
\(\mathcal O_Y(-A)\), not merely a function with the same divisor.  The
chosen generators also show that this line bundle is relatively generated
and hence relatively nef.  It is therefore one of the semipositive model
generators of the revised CLD psh sheaf; no inference from positivity of the
older current predicate is involved.  After pullback to an analytic curve,
the sections generate a semipositive model line bundle on a semistable curve
model.  The outgoing-slope sum at a vertex is its nonnegative degree there,
proving subharmonicity.  Finally,
\eqref{eq:frontier-local-ma-gauss-ideal} identifies the Gauss restriction
with \(-v(\mathfrak a)\), the function of the nef Cartier b-divisor \(-A\).
Tensor powers, multilinearity, and an algebraizable common model give the
asserted rational \(\mathbf Q\)-Cartier case.  No comparison is claimed here
for a model carried only by a nonalgebraizable analytic modification.
\end{proof}

\begin{theorem}\label{thm:frontier-local-ma-model-comparison}
Let \(\mathfrak a_1,\ldots,\mathfrak a_d\) be
\(\mathfrak m\)-primary ideals.  On one log resolution write
\[
 \mathfrak m\mathcal O_Y=\mathcal O_Y\!\left(-\sum_Eb_EE\right),
 \qquad
 \mathfrak a_i\mathcal O_Y=\mathcal O_Y(-A_i).
\]
Then
\begin{equation}
 \MA^{\partial}_{\mathrm{CLD},\mathfrak m}
       (u_{\mathfrak a_1},\ldots,u_{\mathfrak a_d})
 =\sum_E b_E\bigl((-A_1)\cdots(-A_d)\cdot E\bigr)
       \delta_{\iota(v_E)}.
 \label{eq:frontier-local-ma-cld-atomic}
\end{equation}
After applying \(p\), this is exactly
\begin{equation}
 \MA_{\mathrm{BFJ},\mathfrak m}
       (g_{\mathfrak a_1},\ldots,g_{\mathfrak a_d})
 =\sum_E b_E\bigl((-A_1)\cdots(-A_d)\cdot E\bigr)
       \delta_{v_E}.
 \label{eq:frontier-local-ma-bfj-atomic}
\end{equation}
The same comparison holds for rational semipositive model functions on an
algebraizable common log model.
\end{theorem}

\begin{proof}
By \Cref{lem:frontier-local-ma-common-model-core}, the ideal data may be
carried by an algebraic simultaneous log resolution
\(\pi:Y\to U\), where \(U\subset\mathbf A^n_{\mathbf C}\) is a Zariski
neighborhood of the origin and \(\pi\) is an isomorphism over
\(U\setminus\{0\}\).  Compactify \(Y\) and the finitely many divisors without changing a
neighborhood of the proper exceptional fibre.  Let
\(\mathscr Y_0\) be the \(t\)-adic completion of
\(Y\times_{\mathbf C}\Spec K^\circ\), set
\(\mathscr J=(t,\mathcal O_Y(-D))\), and put
\begin{equation}
 q:\mathscr Y:=
 \operatorname{Nor}\!\left(
       \operatorname{Bl}_{\mathscr J}(\mathscr Y_0)\right)
       \longrightarrow\mathscr Y_0\longrightarrow Y,
 \qquad D=\sum_Eb_EE.
 \label{eq:frontier-local-ma-formal-blowup}
\end{equation}
Take a toroidal resolution of this formal model.  The ideal \(\mathscr J\)
contains \(t\), and every subsequent toroidal centre lies in the special
fibre.  Thus all these modifications are admissible, their generic fibres
are unchanged, and the resulting flat formal model has pure relative
dimension \(n\).  Moreover, \(\pi_K^{\An}\) is an isomorphism over the
punctured collar of \(S_{\mathfrak m}\).  The products there are therefore
literally the pullbacks of the products on \(B\); no comparison or
birational-invariance theorem is being assumed.  Compactification is used
only to make the special-fibre components proper, as required by the formal
intersection formula below.  Every component that survives the swept
difference lies over the original proper exceptional fibre, so its
coefficient is independent of this compactification.

Let \(\mathscr F\) be the Cartier divisor determined by
\[
 (t,\mathcal O_Y(-D))\mathcal O_{\mathscr Y}
       =\mathcal O_{\mathscr Y}(-\mathscr F).
\]
The line bundles producing the \(n\) weights in the first term of
\eqref{eq:frontier-local-ma-swept-product} are
\begin{equation}
 \mathcal L_i=q^*\mathcal O_Y(-A_i)\quad(1\leq i\leq d),
 \qquad
 \mathcal L_\chi=\mathcal O_{\mathscr Y}(-\mathscr F).
 \label{eq:frontier-local-ma-model-line-bundles}
\end{equation}
Indeed, the two sections \(t,s_D\) give the weight
\(\log\max\{|t|,|s_D|\}=\chi_{\mathfrak m}\), while the ideal generators
give the first \(d\) bundles.  More explicitly, if \(\tau\) locally
generates \(\mathscr J\mathcal O_{\mathscr Y}\), write
\(t=\tau a\) and \(s_D=\tau b\).  The pair \((a,b)\) generates the unit
ideal, hence \(\max\{|a|,|b|\}=1\) on the corresponding formal domain and
\[
             \log\max\{|t|,|s_D|\}=\log|\tau|.
\]
Thus \(\mathcal L_\chi\) is exactly the tautological model bundle.  The
factorizations \(f_{ij}=s_{A_i}h_{ij}\) and
\(\max_j|h_{ij}|=1\), proved in
\Cref{lem:frontier-local-ma-common-model-core}, similarly identify the
other model bundles.  This verifies the model metrics, not just their
Cartier divisors.

We use the formal-model measure formula
\cite[Theorem~18.5.5 and Proposition~18.5.8]{CLD25}.  In the present
discretely valued, algebraically closed residue-field setting, if \(G\) is a
special-fibre component and
\(m_G=\ord_G(t)\), it reads
\begin{equation}
 c_1(\mathcal M_1)\wedge\cdots\wedge c_1(\mathcal M_n)
 =\sum_Gm_G\deg_G(\mathcal M_1\cdots\mathcal M_n)\,\delta_{\xi_G}.
 \label{eq:frontier-local-ma-model-measure}
\end{equation}
This is also obtained by multilinearity and successive model
Poincar\'e--Lelong intersections.  The normalization of the analytic point is
\(\ord_G/\ord_G(t)\), explaining the factor \(m_G\).

At the generic point of a prime \(E\subset D\), write a local equation of
\(D\) as \(x^{b_E}\).  The normalized blowup of \((t,x^{b_E})\) has one
boundary component \(\mathscr E_E\) with
\begin{equation}
 \ord_{\mathscr E_E}(t)=b_E,
 \qquad
 \ord_{\mathscr E_E}(x)=1.
 \label{eq:frontier-local-ma-boundary-orders}
\end{equation}
Here is the local calculation behind both assertions.  Over the function
field of \(E\), the Rees algebra has the presentation
\[
 \mathbf C(E)[x,t][T_0,T_1]/(x^{b_E}T_0-tT_1).
\]
Its exceptional fibre is
\(\Proj\mathbf C(E)[T_0,T_1]=\mathbf P^1_{\mathbf C(E)}\).
The two affine charts have coordinate rings
\(\mathbf C(E)[x,u]\), with \(t=x^{b_E}u\), and
\(\mathbf C(E)[x,t,v]/(x^{b_E}-tv)\).  The second is smooth for \(b_E=1\);
otherwise it is a hypersurface whose singular locus has codimension two,
hence is normal.  The first is smooth.
They show that the exceptional valuation is the primitive monomial valuation
in \eqref{eq:frontier-local-ma-boundary-orders}, and that normalization does
not change the generic exceptional curve.  Thus the reduced component
\(\mathscr E_E\) maps properly to \(E\) and is generically a
\(\mathbf P^1\)-bundle.  Although
\(\mathscr F\) has coefficient \(b_E\) along \(\mathscr E_E\), its
invertible ideal is the pullback of the Rees tautological
\(\mathcal O(1)\).  Consequently
\(\mathcal L_\chi|_{p_E^{-1}(\eta_E)}\simeq
\mathcal O_{\mathbf P^1}(1)\), so its degree is one, not \(b_E\).  The
other bundles are pulled back from \(Y\), and hence
\begin{equation}
 \mathcal L_i|_{\mathscr E_E}
   =p_E^*\bigl(\mathcal O_Y(-A_i)|_E\bigr).
 \label{eq:frontier-local-ma-boundary-restrictions}
\end{equation}
The projective-bundle pushforward formula therefore gives
\begin{equation}
 \deg_{\mathscr E_E}
       (\mathcal L_1\cdots\mathcal L_d\mathcal L_\chi)
       =(-A_1)\cdots(-A_d)\cdot E.
 \label{eq:frontier-local-ma-boundary-degree}
\end{equation}
Equations \eqref{eq:frontier-local-ma-model-measure},
\eqref{eq:frontier-local-ma-boundary-orders}, and
\eqref{eq:frontier-local-ma-boundary-degree} yield the \(E\)-summand in
\eqref{eq:frontier-local-ma-cld-atomic} at the point
\(\iota(v_E)\).

It remains to account for every other component.  Where
\(u_{\mathfrak m}>-1\), the section \(s_D\) dominates \(t\), so
\(\mathcal L_\chi\) agrees with the model bundle of
\(u_{\mathfrak m}\); the second term of
\eqref{eq:frontier-local-ma-swept-product} cancels its coefficient.  Where
\(u_{\mathfrak m}<-1\), the section \(t\) dominates and
\(\mathcal L_\chi\) is trivial, so the coefficient vanishes.  Over the
generic point of a divisor \(E\), the preceding two-generator Rees
calculation shows that \(\mathscr E_E\) is the unique component on the
equality locus; every other ray lies strictly on one of these two sides.

It remains only to exclude equality-locus components above
higher-codimension strata.  Let \(G\) be such a component, let
\(Z=q(G)\), and put
\[
 \alpha_G:=c_1(\mathcal L_\chi|_G)\cap[G]
             \in A_{n-1}(G).
\]
If \(\operatorname{codim}_Y Z\geq2\), then \(\dim Z\leq n-2\).  By the definition of
proper pushforward of cycles, \(q_*\alpha_G=0\) in \(A_{n-1}(Y)\).  Since
the remaining \(d=n-1\) line bundles are pulled back from \(Y\), the
intersection-theoretic projection formula gives
\[
 \deg_G\!\left(\mathcal L_1\cdots\mathcal L_d
                     \mathcal L_\chi\right)
 =\deg_Y\!\left(c_1(\mathcal O_Y(-A_1))\cdots
        c_1(\mathcal O_Y(-A_d))\cap q_*\alpha_G\right)=0.
\]
Hence the components \(\mathscr E_E\) are the only survivors.  Relative
nefness makes every displayed coefficient nonnegative.

On the BFJ side, \(g_{\mathfrak a_i}\) is the function of the nef Cartier
b-divisor \(-A_i\).  The local model formula
\cite[Proposition~4.9]{BFJ08}, reproduced as
\eqref{eq:local-atomic-ma}, assigns to \(v_E\) the coefficient in
\eqref{eq:frontier-local-ma-bfj-atomic}.  Thus the two formulas agree term by
term; no comparison theorem is imported.  The proof used only multilinear
intersection calculations, so it also applies to rational semipositive model
functions on an algebraizable common log model.

Finally, the common total mass is
\begin{equation}
 \sum_Eb_E\bigl((-A_1)\cdots(-A_d)\cdot E\bigr)
 =(-A_1)\cdots(-A_d)\cdot D,
 \label{eq:frontier-local-ma-total-mass}
\end{equation}
which is the mixed Hilbert--Samuel multiplicity with the fixed reference
ideal \(\mathfrak m\).  This follows by taking general sections of the
\(\mathfrak a_i\) and \(\mathfrak m\) and applying the projection formula.
In particular, all \(\mathfrak a_i=\mathfrak m\) give mass
\(e(\mathfrak m)=1\).  There is no division by
\eqref{eq:frontier-local-ma-total-mass}.
\end{proof}

\begin{proposition}\label{prop:frontier-local-ma-independence}
The common model measure of
\Cref{thm:frontier-local-ma-model-comparison} is independent of the log
resolution, formal model, and tropical charts.  With the boundary datum
\((\mathfrak m,-1)\) fixed, it is local on \(V_{\mathfrak m}\).  Literal
restriction to a relatively open subset is restriction of the measure, with
no renormalization of its mass.
\end{proposition}

\begin{proof}
Two resolutions and two formal models admit a common domination.  If
\(q:Y'\to Y\) is such a domination and \(E'\) is the strict transform of
\(E\), the projection formula gives
\[
 ((-q^*A_1)\cdots(-q^*A_d)\cdot E')
   =((-A_1)\cdots(-A_d)\cdot E).
\]
For a new \(q\)-exceptional prime \(F\), one has \(q_*F=0\) as a divisor
cycle, and the same projection formula gives zero intersection with the
pulled-back factors.  Thus refinement adds no mass and preserves every old
coefficient.  The normalization of the analytic point is also preserved.

Two tropical charts admit a common polyhedral refinement.  On it the
weights are the same intrinsic model functions and the calculation in
\eqref{eq:frontier-local-ma-model-line-bundles} is unchanged.  Finally, both
terms in \eqref{eq:frontier-local-ma-swept-product} are local
Bedford--Taylor measures, and their difference is shell-supported.  Model
potentials coming from the germ factor through \(p\) on the shell, so
agreement near an open subset of the link gives agreement near its inverse
image.  Pushing forward by \(p\) proves locality.  Restriction of a measure
already has this property; changing its total mass would destroy it.
\end{proof}

Define the joint bounded closure \(\mathcal C_b^{\mathrm{com}}\) to consist
of bounded continuous BFJ formal psh functions \(g\) for which there are
decreasing semipositive model functions \(g^{(k)}\downarrow g\), each in
the preceding ideal/algebraizable model locus, whose homogeneous Gauss
extensions \(u^{(k)}\), on one fixed collar of the shell, satisfy all three
conditions:
\begin{enumerate}
\item they belong to the revised CLD psh sheaf;
\item they decrease locally boundedly to a continuous revised-CLD psh
      function \(u\);
\item their restrictions to all analytic curves are subharmonic.
\end{enumerate}
For a finite collection we require a common collar and synchronized common
algebraizable model refinements.  For a further decreasing sequence,
``joint'' means that one decreasing two-index array works simultaneously for
the outer and model indices.

\begin{proposition}\label{prop:frontier-local-ma-bounded-closure}
For \(g_1,\ldots,g_d\in\mathcal C_b^{\mathrm{com}}\), the limit
\begin{equation}
 \MA^{\mathrm{com}}_{\mathfrak m}(g_1,\ldots,g_d)
 :=\lim_{k\to\infty}
 p_*\MA^{\partial}_{\mathrm{CLD},\mathfrak m}
       (u_1^{(k)},\ldots,u_d^{(k)})
 \label{eq:frontier-local-ma-bounded-limit}
\end{equation}
exists, is independent of all approximants, and equals the BFJ mixed local
intersection measure.  It is continuous along decreasing sequences admitting
one joint bounded model approximation of the preceding kind.
\end{proposition}

\begin{proof}
At every model index,
\Cref{thm:frontier-local-ma-model-comparison} gives
\begin{equation}
 p_*\MA^{\partial}_{\mathrm{CLD},\mathfrak m}
       (u_1^{(k)},\ldots,u_d^{(k)})
 =\MA_{\mathrm{BFJ},\mathfrak m}
       (g_1^{(k)},\ldots,g_d^{(k)}).
 \label{eq:frontier-local-ma-model-index-equality}
\end{equation}
Let \(B\) denote the symmetric \(d\)-linear BFJ measure operator.  Its
polarization identity is
\begin{equation}
 d!B(h_1,\ldots,h_d)
 =\sum_{J\subset\{1,\ldots,d\}}(-1)^{d-|J|}
 B\!\left(\sum_{j\in J}h_j,\ldots,
          \sum_{j\in J}h_j\right).
 \label{eq:frontier-local-ma-polarization}
\end{equation}
For every \(J\), the diagonal input on the right decreases to the
corresponding sum of the limiting functions.  BFJ bounded monotone
continuity \cite[Proposition~4.8]{BFJ08}, followed by the finite sum in
\eqref{eq:frontier-local-ma-polarization}, proves weak convergence of the
mixed measures.  Together with
\eqref{eq:frontier-local-ma-model-index-equality}, this proves existence of
\eqref{eq:frontier-local-ma-bounded-limit}.  Notice that the argument uses
the swept model measures; it never commutes a characteristic function with a
weak limit.

The limiting BFJ product depends only on \(g_1,\ldots,g_d\), hence proves
independence.  For a further decreasing sequence, the joint approximation
hypothesis supplies an array decreasing in both indices.  Its diagonal is a
cofinal decreasing sequence converging to the final limit, so the same
polarization argument proves the asserted continuity.  On curves, the
submean, equivalently outgoing-slope, inequalities pass to the continuous
decreasing limit.  No claim is made for a bounded curve-test psh function
without this joint revised-CLD/BFJ approximation.
\end{proof}

\begin{proposition}\label{prop:frontier-local-ma-reference-ideal}
Let \(I\) be an \(\mathfrak m\)-primary ideal.  For ideal functions, and for
rational semipositive functions carried by an algebraizable common log
model, the preceding construction and comparison remain valid with the link
\(V_I=\{v:v(I)=1\}\), the shell \(u_I=-1\), and the cutoff
\(\max\{u_I,-1\}\).  The joint bounded extension remains valid with the
same approximation hypothesis.  We denote the resulting measure by
\(\MA_I^{\mathrm{com}}\).
\end{proposition}

\begin{proof}
Polynomialize \(I\) and the finitely many input ideals simultaneously by the
Nakayama argument in
\Cref{lem:frontier-local-ma-common-model-core}, and principalize them on one
algebraic log model.  Write
\(I\mathcal O_Y=\mathcal O_Y(-D_I)\).  The Gauss extension in
\eqref{eq:frontier-local-ma-gauss-extension} uses only the normalization
\(v(I)=\widetilde v(t)=1\), while primaryness bounds all input potentials on
a fixed collar of the shell.  Replace \(\mathscr J\) in
\eqref{eq:frontier-local-ma-formal-blowup} by
\((t,\mathcal O_Y(-D_I))\).  At the generic point of a prime \(E\subset
D_I\), the same Rees calculation is the normalized blowup of
\((t,x^{\ord_E(I)})\).  It gives the coefficient
\[
 \ord_E(I)\bigl((-A_1)\cdots(-A_d)\cdot E\bigr),
\]
which is precisely the \(I\)-normalized BFJ coefficient.  The cancellation
off the shell and the higher-codimension projection argument are unchanged.
Finally, the proof of
\Cref{prop:frontier-local-ma-bounded-closure} uses only this model equality
and bounded monotone convergence, so it applies verbatim under the stated
joint approximation hypothesis.
\end{proof}

\begin{proposition}\label{prop:frontier-local-ma-finite-map}
Let \(F:(X',0')\to(X,0)\) be a finite map of smooth
\(n\)-dimensional germs of local degree \(q\), and put
\(I=\mathfrak m\mathcal O_{X',0'}\).  Normalize the source link first by
\(I\), and write
\(F_\#:V_I\to V_{\mathfrak m}\), \(w\mapsto w\circ F^\#\).  For ideal
model functions, for rational functions whose target models and pulled-back
source models both lie in the algebraizable model locus, and for members of the joint bounded
closure whose pullbacks remain in that closure,
\begin{equation}
 (F_\#)_*\MA^{\mathrm{com}}_I(F^*g_1,\ldots,F^*g_d)
 =q\,\MA^{\mathrm{com}}_{\mathfrak m}(g_1,\ldots,g_d).
 \label{eq:frontier-local-ma-finite-projection}
\end{equation}
On the standard source link \(V_{\mathfrak m'}\), set
\[
 c(v):=v(I),\qquad R_I(v):=v/c(v),\qquad
 \widehat F(v):=(v\circ F^\#)/v(I)=F_\#(R_I(v)).
\]
For ideal model functions, the exact standard-link formula is
\begin{equation}
 \widehat F_*\!\left(
 c\,\MA^{\mathrm{com}}_{\mathfrak m'}
   (g_{F^*\mathfrak a_1},\ldots,g_{F^*\mathfrak a_d})\right)
 =q\,\MA^{\mathrm{com}}_{\mathfrak m}
   (g_{\mathfrak a_1},\ldots,g_{\mathfrak a_d}).
 \label{eq:frontier-local-ma-standard-link}
\end{equation}
For an \'etale map of pointed complex germs, \(q=1\),
\(I=\mathfrak m'\), and \(c\equiv1\).
\end{proposition}

\begin{proof}
Choose compatible analytic log resolutions, and let \(h:Y'\to Y\) be the
normalization of the analytic fibre product before resolving the source.
The normalization is finite of degree \(q\).  Divisors created
only by the final resolution map to subsets of codimension at least two and
have zero coefficient by the projection formula.  Let \(E'\) be a prime on
the normalized source model that dominates a target prime \(E\).  Denote its
ramification and residue degrees by \(e(E'/E)\) and \(f(E'/E)\).  Then
\begin{equation}
 \ord_{E'}(I)=e(E'/E)b_E,
 \qquad h_*E'=f(E'/E)E,
 \label{eq:frontier-local-ma-finite-ef}
\end{equation}
and, if
\(\alpha_E=(-A_1)\cdots(-A_d)\cdot E\), the ordinary projection formula
gives
\begin{equation}
 \bigl((-h^*A_1)\cdots(-h^*A_d)\cdot E'\bigr)
       =f(E'/E)\alpha_E.
 \label{eq:frontier-local-ma-finite-intersection}
\end{equation}
Consequently the source atom over \(E'\) has coefficient
\(e(E'/E)f(E'/E)b_E\alpha_E\), and its image under \(F_\#\) is
\(v_E\).  The equality
\(\sum_{E'\mid E}e(E'/E)f(E'/E)=q\) now gives exactly the coefficient
\(qb_E\alpha_E\).  This proves
\eqref{eq:frontier-local-ma-finite-projection} term by term for ideal model
functions.  Multilinearity proves the algebraizable rational model case;
here the hypothesis on the pullback is needed because the finite analytic map
itself need not be algebraic.  For ideal functions, the source and target
CLD--BFJ comparisons may be computed on their respective algebraic models;
model independence identifies both with the analytic intersection
calculation above.  The same reasoning applies to the stated algebraizable
rational locus.  Joint bounded monotone approximation and continuity of
pushforward of finite measures prove the bounded extension.

To change the reference ideal, let
\(v_{E'}=\ord_{E'}/b'_{E'}\in V_{\mathfrak m'}\) and put
\(c_{E'}=\ord_{E'}(I)\).  Then
\[
 c(v_{E'})=c_{E'}/b'_{E'},
 \qquad
 R_I(v_{E'})=\ord_{E'}/c_{E'}.
\]
Put
\(\beta_{E'}=(-A'_1)\cdots(-A'_d)\cdot E'\).  The coefficient of this atom
for the \(\mathfrak m'\)-reference is \(b'_{E'}\beta_{E'}\), whereas its
coefficient for the \(I\)-reference is \(c_{E'}\beta_{E'}\).  Multiplication
by \(c(v_{E'})=c_{E'}/b'_{E'}\) therefore gives the latter coefficient
exactly, and radial normalization changes only the location of the atom.
For clarity, write
\(g_{\mathfrak b}^{[J]}(v):=-v(\mathfrak b)\) on \(V_J\).  Hence, for the
pulled-back ideals \(\mathfrak b_i=\mathfrak a_i\mathcal O_{X',0'}\),
\begin{equation}
 \MA^{\mathrm{com}}_I
   (g_{\mathfrak b_1}^{[I]},\ldots,g_{\mathfrak b_d}^{[I]})
 =R_{I*}\!\left(c\,\MA^{\mathrm{com}}_{\mathfrak m'}
   (g_{\mathfrak b_1}^{[\mathfrak m']},\ldots,
    g_{\mathfrak b_d}^{[\mathfrak m']})\right).
 \label{eq:frontier-local-ma-reference-change}
\end{equation}
Combining \eqref{eq:frontier-local-ma-finite-projection} and
\eqref{eq:frontier-local-ma-reference-change} gives
\eqref{eq:frontier-local-ma-standard-link}.  The weight is genuinely
variable: for
\(F(z_1,\ldots,z_n)=(z_1^q,z_2,\ldots,z_n)\),
\[
 v(\mathfrak m\mathcal O_{X',0'})
   =\min\{qv(z_1),v(z_2),\ldots,v(z_n)\}
\]
is not constant on \(V_{\mathfrak m'}\).
\end{proof}

\begin{proposition}\label{prop:frontier-local-ma-no-probability}
No rule can both replace every nonzero literal local restriction of the
common measure by a probability measure and form a sheaf under literal
restriction.
\end{proposition}

\begin{proof}
Work on \((\mathbf C^2,0)\).  Blow up the origin and then the point of the
first exceptional curve corresponding to the tangent direction \(y=0\).
Let \(E_1\) be the strict transform of the first exceptional curve and
\(E_2\) the second.  Then
\[
 E_1^2=-2,\qquad E_2^2=-1,\qquad E_1\cdot E_2=1,
 \qquad b_{E_1}=b_{E_2}=1.
\]
For \(\mathfrak a=(x^3,xy,y^2)\), the divisorial orders are two and three,
so \(A=2E_1+3E_2\).  Direct intersection gives
\[
 (-A)\cdot E_1=-\bigl(2(-2)+3\bigr)=1,
 \qquad
 (-A)\cdot E_2=-\bigl(2+3(-1)\bigr)=1.
\]
Thus \eqref{eq:frontier-local-ma-bfj-atomic} gives
\(\mu=\delta_{v_{E_1}}+\delta_{v_{E_2}}\).  Choose an open set \(U\)
containing both atoms and \(W\subset U\) containing only \(v_{E_1}\).
Probability normalization on \(U\), followed by restriction to \(W\), gives
\(\tfrac12\delta_{v_{E_1}}\), whereas direct probability normalization on
\(W\) gives \(\delta_{v_{E_1}}\).  They differ.  The unnormalized common
measure itself restricts correctly.
\end{proof}

The next example shows why the bounded hypothesis cannot simply be deleted.

\begin{example}\label{ex:frontier-local-ma-polar-endpoint}
On \((\mathbf C^2,0)\), let
\(v_s(x)=1\), \(v_s(y)=s\) for \(1\leq s<+\infty\), and let
\(v_\infty\) be the endpoint with \(v_\infty(y)=+\infty\).  For
\(\mathfrak a_k=(y,x^k)\),
\begin{equation}
 g_k(v_s)=-v_s(\mathfrak a_k)=-\min\{s,k\}\downarrow-s=:g(v_s),
 \label{eq:frontier-local-ma-unbounded-ray}
\end{equation}
and \(g(v_\infty)=-\infty\).  Resolving \(\mathfrak a_k\) gives a chain with
\(v_j(x)=1\), \(v_j(y)=j\), exceptional intersections
\(E_j^2=-2\) for \(j<k\), \(E_k^2=-1\), and adjacent intersection one.
The ideal divisor has coefficient \(j\) along \(E_j\).  Hence
\((-A)\cdot E_j=0\) for \(j<k\) and
\((-A)\cdot E_k=1\), so
\begin{equation}
                  \MA_{\mathrm{BFJ}}(g_k)=\delta_{v_k}.
 \label{eq:frontier-local-ma-unbounded-atoms}
\end{equation}
The full decreasing limit is \(\delta_{v_\infty}\).  On the other hand,
\(g_k=\max\{g,-k\}\), and the nonpluripolar truncation
\[
 \one_{\{g>-k\}}\MA_{\mathrm{BFJ}}(\max\{g,-k\})
\]
is zero: the only atom in
\eqref{eq:frontier-local-ma-unbounded-atoms} lies on \(\{g=-k\}\).  Thus
full monotone continuation and nonpluripolar continuation disagree by a
polar endpoint mass.
\end{example}

\begin{theorem}\label{thm:frontier-common-local-ma}
For a smooth complex germ and the fixed datum
\((\mathfrak m,|t|=\rme^{-1},u_{\mathfrak m}=-1)\), the Gauss realization
of \Cref{lem:frontier-local-ma-gauss} identifies the swept CLD boundary
measure with the BFJ local measure on ideal functions and on rational
semipositive functions carried by an algebraizable common log model.  This
identification is model- and chart-independent within that locus, local
without probability renormalization, and functorial under finite maps in the
input domains specified in \Cref{prop:frontier-local-ma-finite-map}, with the
reference-change factor in \eqref{eq:frontier-local-ma-reference-change}.
It extends uniquely to
\(\mathcal C_b^{\mathrm{com}}\) with joint decreasing monotone continuity.
Without an additional convention governing polar mass, the model comparison
does not determine an extension to arbitrary unbounded functions.
\end{theorem}

\begin{proof}
The actual analytic/polyhedral realization is
\Cref{lem:frontier-local-ma-gauss}; simultaneous membership in the three
model positivity classes is
\Cref{lem:frontier-local-ma-common-model-core}; and the term-by-term
intersection equality is
\Cref{thm:frontier-local-ma-model-comparison}.  Independence and locality
are \Cref{prop:frontier-local-ma-independence}, the bounded extension is
\Cref{prop:frontier-local-ma-bounded-closure}, and the finite-map formulas
are \Cref{prop:frontier-local-ma-finite-map}.  The incompatibility of
probability normalization with restriction is
\Cref{prop:frontier-local-ma-no-probability}.  Finally,
\Cref{ex:frontier-local-ma-polar-endpoint} gives a full decreasing limit that
charges the polar endpoint, whereas nonpluripolar truncation discards that
mass.  Thus the model comparison alone does not select the treatment of polar
mass.  These results prove all assertions and the final limitation.
\end{proof}

Consequently the model-locus comparison and the explicitly defined joint
bounded closure are \emph{verified}.  The following parts of the problem
remain \emph{open}: a product for an arbitrary
bounded classically psh function not known to admit the joint realization; a
capacity or finite-energy criterion selecting the treatment of polar mass;
comparison for a rational model carried only by a nonalgebraizable analytic
modification; and extensions to singular or nonnormal germs.  No primary
source located by the cutoff states the comparison theorem above; its proof
is a new Rethlas-verified manuscript argument, not a claim about priority in
the literature.  The Bingener comparison recorded in
\Cref{thm:mp-bingener-projective-germs} and based on
\cite[Aussage~2.8]{Bingener76Stein} algebraizes suitable projective germs over
a local or Stein algebra; it does not by itself produce a finite-type
\(\mathbf C\)-model and an admissible CLD formal model for an arbitrary
analytic modification, so it does not close the stated algebraizability gap.

\subsection{Regularity and orthogonality on balanced polyhedral spaces}

\begin{problem}\label{prob:frontier-polyhedral-regularity}
Let \(\mathbf X\) be a pure-dimensional balanced rational polyhedral space
and let \(\gamma\) be a piecewise-affine polyhedrally convex reference
function positive at infinity.  Characterize the pairs
\((\mathbf X,\gamma)\) for which the
regularity and orthogonality properties of Botero--Mazzon--Pille-Schneider
hold.  Under realizability hypotheses, compare their polyhedral
Monge--Amp\`ere operator with CLD, delta-form, and non-Archimedean
Monge--Amp\`ere measures.
\end{problem}

The bibliographic dates require a correction.  The first arXiv version was
submitted on 9 March 2026 and the date printed inside the manuscript is
10 March 2026, not 11 August 2026 \cite{BMPS26}.  We use that version
throughout.

Here are the conventions needed to compare it with
\Cref{app:tropical-algebra}.  Normal vectors point from a codimension-one
face into an adjacent top-dimensional face, as in
\eqref{eq:appG-primitive-normal}, and a convex bend has positive divisor.
Botero--Mazzon--Pille-Schneider allow positive real balancing weights; a
tropical cycle in the narrower Mikhalkin--Rau convention has integral
weights.  Analytic comparison therefore uses the latter only after the
saturated tangent lattices and analytic tropical multiplicities have been
installed.  Finally, their real Monge--Amp\`ere measure is the Alexandrov
measure of \Cref{thm:appD-alexandrov-measure}, without a factorial.  Their
polyhedral intersection convention consequently gives
\begin{equation}
 \one_{\RelInt(\sigma)}\MA_{\mathrm{poly}}(\varphi)
 =d![\mathbf X](\sigma)\,
   \MA_{\mathbf R}(\varphi|_{\RelInt(\sigma)})
 \label{eq:frontier-polyhedral-real-local}
\end{equation}
on the relative interior of every top-dimensional face
\cite[Propositions~4.12 and~4.20]{BMPS26}.  This is exactly the factorial
firewall in \Cref{app:normalizations}.

We first separate the source results from the new deduction below.  For a
piecewise-affine polyhedrally convex reference \(\gamma\) positive at
infinity, \cite[Theorem~4.13 and Proposition~4.16]{BMPS26} constructs the
positive mixed operator
\[
 \MA_{\mathrm{poly}}:
 \mathrm{PC}_{\mathrm{reg}}(\mathbf X,\gamma)^d
 \longrightarrow\mathcal M^+(\mathbf X),
 \qquad
 \int_{\mathbf X}\MA_{\mathrm{poly}}=\deg(\gamma),
\]
and proves its continuity under uniform, in particular decreasing,
regularizable approximation.  Its Theorem~6.4 says that regularity is
equivalent to each of
\[
 \begin{split}
 &P_\gamma(u)\in\mathrm{PC}_{\mathrm{reg}}(\mathbf X,\gamma)
       \quad\text{for every bounded piecewise-affine }u,\\
 &\mathrm{PAPC}(\mathbf X,\gamma)
       \text{ is uniformly dense in }\mathrm{PC}(\mathbf X,\gamma),\\
 &\mathrm{PC}(\mathbf X,\gamma)
       =\mathrm{PC}_{\mathrm{reg}}(\mathbf X,\gamma).
 \end{split}
\]
Regularity holds in dimension one, on the complete affine space
\(N_{\mathbf R}\), and for conical functions in the conical category
\cite[Proposition~6.13 and Remark~6.6]{BMPS26}.  Orthogonality is the further
identity
\begin{equation}
 \int_{\mathbf X}
 \bigl(P_\gamma(u)-(\gamma+u)\bigr)
 \MA_{\mathrm{poly}}(P_\gamma(u))=0.
 \label{eq:frontier-polyhedral-orthogonality}
\end{equation}
The source proves \eqref{eq:frontier-polyhedral-orthogonality} on
one-dimensional polyhedrally smooth spaces and obtains existence and
uniqueness, modulo constants, for compactly supported positive measures of
mass \(\deg(\gamma)\)
\cite[Proposition~6.21 and Theorem~6.22]{BMPS26}.

The following elementary convex-envelope fact is the missing contact
mechanism on a complete affine space.

% RETHLAS-VERIFIED: balanced-polyhedral-affine-contact; receipt 20260814T174427Z_438bf42ae745
\begin{lemma}\label{lem:frontier-polyhedral-affine-contact}
Let \(f:\mathbf R^d\to\mathbf R\) be continuous and coercive, and let
\(P\) be its greatest lower-semicontinuous convex minorant.  Suppose that
\(P\) is polyhedral.  If \(\partial P(x)\) has positive
\(d\)-dimensional volume, then \(P(x)=f(x)\).
\end{lemma}

\begin{proof}
Choose a Euclidean norm and a point \(p\) in the interior of
\(\partial P(x)\).  For some \(\eta>0\),
\(p+\eta\overline B^\vee\subset\partial P(x)\).  For \(y\ne x\), choose
\(q_y\) in the dual unit sphere with
\(\langle q_y,y-x\rangle=\lVert y-x\rVert\).  Applying the subgradient
inequality to \(p+\eta q_y\) gives
\begin{equation}
 P(y)-\langle p,y\rangle
 \geq P(x)-\langle p,x\rangle+\eta\lVert y-x\rVert.
 \label{eq:frontier-polyhedral-exposed-contact}
\end{equation}
Thus \(x\) is the unique minimizer of
\(P-\langle p,\cdot\rangle\).

Set
\[
                 m:=\inf_y\{f(y)-\langle p,y\rangle\}.
\]
The affine function \(y\mapsto\langle p,y\rangle+m\) is bounded above by
\(f\); maximality of \(P\) therefore puts it below \(P\).  Since
\(P\leq f\), the two opposite inequalities between the infima give
\begin{equation}
 m=\inf_y\{P(y)-\langle p,y\rangle\}
  =P(x)-\langle p,x\rangle.
 \label{eq:frontier-polyhedral-contact-infimum}
\end{equation}
By \eqref{eq:frontier-polyhedral-exposed-contact} and \(f\geq P\), the
function \(f-\langle p,\cdot\rangle\) is coercive, so it attains its
infimum at some \(y_0\).  Equations
\eqref{eq:frontier-polyhedral-exposed-contact}--
\eqref{eq:frontier-polyhedral-contact-infimum} force \(y_0=x\), and equality
at the minimum gives \(f(x)=P(x)\).
\end{proof}

% RETHLAS-VERIFIED: balanced-polyhedral-affine-orthogonality; receipt 20260814T174427Z_438bf42ae745
\begin{theorem}\label{thm:frontier-polyhedral-affine-orthogonality}
Let \(N\) have rank \(d\), let
\(\mathbf X=N_{\mathbf R}\) carry the constant balancing weight
\(c>0\), and let \(\gamma\) be a piecewise-affine polyhedrally convex
function positive at infinity.  Then
\((\mathbf X,\gamma)\) has regularity and orthogonality in every dimension.
No strict-convexity hypothesis beyond positivity at infinity is needed.
\end{theorem}

\begin{proof}
Let \(u\) be bounded and piecewise affine, put \(f=\gamma+u\), and choose
\(C\) with \(|u|\leq C\).  Positivity of \(\gamma\) at infinity and
boundedness of \(u\) make \(f\) coercive.  On all of
\(N_{\mathbf R}\), polyhedral
convexity is ordinary convexity.  The envelope \(P_\gamma(u)\) is therefore
the greatest lower-semicontinuous convex minorant \(P\) of \(f\): indeed,
\(\gamma-C\leq f\) is a competitor, while every convex minorant between
\(\gamma-C\) and \(f\) has the required \(\gamma+O(1)\) growth.

On a common finite rational polyhedral decomposition, the epigraph of
\(f\) is a finite union of rational polyhedra.  The closed convex hull of a
finite union of rational polyhedra is a rational polyhedron: this follows,
for example, by writing each member as the convex hull of finitely many
rational points plus the cone generated by finitely many rational recession
directions, and taking the convex hull of the points and the cone generated
by all recession directions.  Since
\[
 \operatorname{epi}(P)
 =\overline{\operatorname{conv}}\bigl(\operatorname{epi}(f)\bigr),
\]
its lower boundary is piecewise affine.  The inequalities
\(
 \gamma-C\leq P\leq f\leq\gamma+C
\)
show that \(P\in\mathrm{PAPC}(\mathbf X,\gamma)\).  This proves the first
condition in the source regularity characterization, hence regularity.

For a piecewise-affine convex function on \(N_{\mathbf R}\), apply the
piecewise-affine version of
\eqref{eq:frontier-polyhedral-real-local} to the coarsest complex with the
single face \(N_{\mathbf R}\); the domains of affinity of \(P\) are merely
an auxiliary refinement.  This gives the global identity
\begin{equation}
                  \MA_{\mathrm{poly}}(P)
                  =d!c\,\MA_{\mathbf R}(P).
 \label{eq:frontier-polyhedral-affine-global-ma}
\end{equation}
The right-hand side is atomic, and its mass at \(x\) is
\(\operatorname{vol}(\partial P(x))\), by
\Cref{ex:appD-polyhedral-atom}.  Every atom of positive mass is therefore
on the contact set \(\{P=f\}\), by
\Cref{lem:frontier-polyhedral-affine-contact}.  The integrand in
\eqref{eq:frontier-polyhedral-orthogonality} vanishes on the support of the
measure, proving orthogonality.
\end{proof}

The regularity assertion in
\Cref{thm:frontier-polyhedral-affine-orthogonality} is the complete case of
\cite[Remark~6.6]{BMPS26}; the all-dimensional orthogonality deduction is
new in this manuscript and has been independently verified by Rethlas.

The next example shows why completeness of the ambient convex geometry
cannot be replaced by strict positivity of the reference alone.

\begin{example}\label{ex:frontier-polyhedral-cross-orthogonality}
Let
\[
 \mathbf X=(\mathbf R\times\{0\})\cup(\{0\}\times\mathbf R)
\]
with all four ray weights equal to one, and put
\(\gamma(x,y)=|x|+|y|\) on \(\mathbf X\).  This reference is positive at
infinity and strictly polyhedrally convex on the four-ray complex.  For
\(\varepsilon>0\), define the obstacle
\[
 f(x,0)=|x|+\varepsilon,
 \qquad
 f(0,y)=\bigl||y|-\varepsilon\bigr|,
 \qquad u=f-\gamma.
\]
Then \(u\) is bounded and piecewise affine, and
\[
 P_\gamma(u)(x,0)=|x|,
 \qquad
 P_\gamma(u)(0,y)=\max\{0,|y|-\varepsilon\}.
\]
\end{example}

Indeed, the displayed candidate is the restriction of the ambient convex
function
\(
 \max\{|x|,|y|-\varepsilon,0\}
\), hence is polyhedrally convex and lies below \(f\).  Conversely, if
\(q\leq f\) is an admissible polyhedrally convex function, convexity between
\((0,\varepsilon)\) and \((0,-\varepsilon)\) gives \(q(0,0)\leq0\).
On either horizontal ray, convexity and the asymptotic slope bound one give
\(q(x,0)\leq|x|\).  On the positive vertical ray, convexity first gives
\(q(0,y)\leq0\) for \(0\leq y\leq\varepsilon\), and then the same slope
bound gives \(q(0,y)\leq y-\varepsilon\) for
\(y\geq\varepsilon\); the negative ray is identical.  This proves the
envelope formula.  Its outgoing slopes at the origin have sum two, so
\[
 \MA_{\mathrm{poly}}(P_\gamma(u))\{(0,0)\}=2,
 \qquad P_\gamma(u)(0,0)-f(0,0)=-\varepsilon.
\]
All other atoms lie on the contact set, and the orthogonality integral is
therefore \(-2\varepsilon\).  Finally, a positive local one-cycle at the
origin has weights \(a,a,b,b>0\) on the two opposing horizontal and two
opposing vertical rays; intersection with \(\gamma\) has coefficient
\(2a+2b>0\).  The strict-convexity criterion of
\cite[Remark~3.9]{BMPS26} applies.  This is the source's Example~6.23 with
its sign and contact calculation made explicit.

There is a complementary smooth-fan warning in
\cite[Example~6.24, Claim~6.25, and Remark~6.26]{BMPS26}.  On the Bergman
fan of the uniform matroid \(U_{3,4}\), assign value zero to the opposite
rays generated by \(e_{12}\) and \(e_{34}=-e_{12}\), and value one to the
other primitive ray generators.  The resulting polyhedrally convex
reference \(\gamma\) admits a bounded piecewise-affine obstacle \(g\geq
\gamma\) for which
\[
 P_\gamma(g-\gamma)=\gamma,
 \qquad
 \MA_{\mathrm{poly}}(\gamma)\{0\}=4,
 \qquad
 \gamma(0)<g(0).
\]
The source proves the envelope identity cone by cone: equality first holds
on the line \(\mathbf R e_{12}\), then on the original rays by one-variable
convexity, and finally on each two-dimensional cone by joining its boundary
to the diagonal ray.  The mass is the direct iterated-divisor computation
on that fan.  This example is tropically smooth but \(\gamma\) is identically
zero on \(\mathbf R e_{12}\).  Hence it is neither strictly convex nor
positive at infinity.  It does not disprove the conjunction of a suitable
higher-dimensional smoothness condition with the strict growth hypotheses
in \Cref{prob:frontier-polyhedral-regularity}.

Balancing also does not, by itself, produce an analytic space.  The exact
formal comparison that is available once a realization has been supplied is
as follows.

% RETHLAS-VERIFIED: balanced-polyhedral-recursive-realization; receipt 20260814T174427Z_438bf42ae745
\begin{proposition}\label{prop:frontier-polyhedral-recursive-realization}
Let \(Y\) be a pure \(d\)-dimensional non-Archimedean analytic space over a
complete field, and let \(h:Y\to\mathbf X\) be a proper harmonic
tropicalization.  Equip \(\mathbf X\) with the saturated integral
lattices and analytic tropical multiplicities induced by \(h\).  Suppose
that every tropical cycle \(c\) obtained from \([\mathbf X]\) by successive
rational piecewise-affine divisors has a CLD/delta current \(T_c\) such
that
\begin{equation}
 T_{[\mathbf X]}=[Y],\qquad
 T_{a\cdot c}=\dprimeddprime[a\circ h]\wedge T_c,
 \qquad h_*T_c=\delta_c.
 \label{eq:frontier-polyhedral-recursive-current}
\end{equation}
Assume also that regularizable polyhedrally convex weights induce
semipositive analytic metrics and that their analytic wedge products are
weakly continuous under the same decreasing uniform piecewise-affine
approximations as \(\MA_{\mathrm{poly}}\).  Then
\begin{equation}
 h_*\bigl(
  \dprimeddprime[a_1\circ h]\wedge\cdots\wedge
  \dprimeddprime[a_d\circ h]\wedge[Y]
 \bigr)
 =\delta_{a_1\cdots a_d\cdot[\mathbf X]}
 =\MA_{\mathrm{poly}}(a_1,\ldots,a_d)
 \label{eq:frontier-polyhedral-realization-comparison}
\end{equation}
for rational piecewise-affine polyhedrally convex weights, and the same
identity holds for regularizable weights.  When the analytic wedge is the
first-Chern, hence Chambert--Loir, measure of the associated semipositive
metrics, it is also the non-Archimedean Monge--Amp\`ere measure.
\end{proposition}

\begin{proof}
Starting with \(T_{[\mathbf X]}=[Y]\), apply the middle identity in
\eqref{eq:frontier-polyhedral-recursive-current} successively.  After
\(j\) steps, the actual analytic wedge is the current attached to the
iterated divisor cycle; hence after \(d\) steps it is
\(T_{a_1\cdots a_d\cdot[\mathbf X]}\).  The last identity in
\eqref{eq:frontier-polyhedral-recursive-current} gives the first equality
in \eqref{eq:frontier-polyhedral-realization-comparison}.  The second is the
definition of the mixed polyhedral Monge--Amp\`ere measure on
piecewise-affine weights.

For regularizable weights, choose decreasing uniform piecewise-affine
approximations in every entry.  Source continuity gives weak convergence on
the polyhedral side, while the stated analytic continuity gives weak
convergence of the wedge currents.  Properness on their supports permits
pushforward through the limit.  This proves the asserted extension.  The
last sentence is simply the identification of the analytic wedge with the
first-Chern/Chambert--Loir measure under the hypothesis stated there.
\end{proof}

The recursive current identity in
\eqref{eq:frontier-polyhedral-recursive-current} is essential.  A pushed
forward one-step Poincar\'e--Lelong formula alone cannot be iterated after the
first divisor, because it has forgotten the current carried by the
intermediate cycle.  Likewise, arbitrary positive balancing weights cannot
replace analytic tropical multiplicities: the latter include the local
degree and saturated lattice index of
\eqref{eq:appG-lattice-index}.  Formula
\eqref{eq:frontier-polyhedral-realization-comparison} commutes with complete
scalar extension only when the realization, calibration, currents, metrics,
and both identities in
\eqref{eq:frontier-polyhedral-recursive-current} do so.

The unconditional analytic result in the source is narrower.  For the
toric degenerations of Calabi--Yau complete intersections in its
\S7, let \(\mathbf X\) be the tropicalization, let
\(\gamma\) be the polarization weight, and let \(\mu_0\) be Lebesgue
measure on the essential skeleton.  Under the three assumptions
\[
 (L^d)=\deg(\gamma),\qquad
 \varphi\in\mathrm{PC}_{\mathrm{reg}}(\mathbf X,\gamma),
 \quad\MA_{\mathrm{poly}}(\varphi)=\mu_0,
\]
and semipositivity of
\(
 \Psi=\varphi_0+(\varphi-\gamma)\circ\operatorname{trop}
\), Proposition~7.2 of \cite{BMPS26} proves
\(
 \MA_{\NA}(\Psi)=\mu_0
\)
and, over \(\mathbf C((t))\), the comparison property.  Its Lemma~7.3
computes every maximal-cell weight as one from analytic degree divided by
lattice index.  The local equality is then
\eqref{eq:frontier-polyhedral-real-local} together with Vilsmeier's
real/non-Archimedean comparison \cite[Theorem~1.1]{Vil20}; equality away
from lower-dimensional skeleton faces and equality of total masses force
the global identity.  None of this is a realization theorem for an abstract
balanced space.

The source itself formulates three sharper tasks that organize this broad
problem: enlarge the notion of polyhedral smoothness so that higher-dimensional
regularity and orthogonality hold; characterize polarized Calabi--Yau
families whose tropicalizations satisfy those properties; and determine
which matroids have Bergman fans satisfying orthogonality
\cite[pp.~3--4, \S1.5, and Remark~6.27]{BMPS26}.  By contrast, analytic
realization of the cross-fan in
\Cref{prob:frontier-cross-fan} is a diagnostic problem formulated in this
book, not a question attributed to Botero--Mazzon--Pille--Schneider.

Thus the audit gives four verdicts.  The complete affine case in
\Cref{thm:frontier-polyhedral-affine-orthogonality} is \emph{verified};
strict positivity alone is \emph{disproved} by
\Cref{ex:frontier-polyhedral-cross-orthogonality}; the general analytic
comparison is \emph{conditional} on the recursive calibrated realization
of \Cref{prop:frontier-polyhedral-recursive-realization}; and the main
higher-dimensional local problem remains \emph{open}.  The missing local
assertion is a contact theorem on singular fans: under an intrinsic
smoothness and strict-positivity package, positive top intersection mass of
an envelope should force contact with the obstacle.  Neither a proof nor a
counterexample satisfying both parts of that package is presently known.

\subsection{Pluriharmonic regularity}

\begin{problem}\label{prob:frontier-pluriharmonic}
Let \(K\) be complete and nontrivially valued and let \(X\) be a good,
separated, strict, boundaryless \(K\)-analytic space.  Is every classically
pluriharmonic function PL without the additional regularity hypotheses in
\cite[Proposition~7.8]{GR25Classical}?  If a continuous \(h\) satisfies
\(\dprimeddprime[h]=0\), is \(h\) classically pluriharmonic, or harmonic under
the hypotheses where a classical harmonic notion is available?  On every
smooth Berkovich curve, does decreasing regularization hold with
CLD-smooth---rather than merely weakly smooth---harmonic approximants?
\end{problem}

Weakly smooth \(\dprimeddprime\)-harmonic functions are harmonic PL by
\cite[Proposition~10.17]{GJR26}.  Distributional vanishing on cell interiors
only proves affinity there; it does not supply the face-balancing law or a
finite harmonic tropicalization.  On curves, weakly smooth and CLD-smooth
regularity are genuinely different.  The search therefore treats the three
questions separately.  The weakly smooth statement is \emph{verified}; the
unrestricted PL, representative, and CLD-smooth claims are \emph{open}.

\subsection{Early collapse of the CLD psh hierarchy}
\index{plurisubharmonic function!CLD hierarchy}
\index{CLD hierarchy!early collapse}

\begin{problem}\label{prob:frontier-cld-hierarchy}
Does the transfinite hierarchy used to generate the CLD sheaf of psh
functions collapse after its first decreasing-limit stage?  If not, is there
a finite or otherwise uniformly controlled ordinal at which it collapses on
every affine tropical space in a natural geometric class?
\end{problem}

This is not the set-theoretic question of eventual stabilization.  The
cardinality argument in \cite[Lemma~8.1.4]{CLD25} already proves that the
hierarchy eventually stabilizes.  Read literally, the sentence in
\cite[Remark~8.1.6]{CLD25} saying that stationarity is not known conflicts
with that lemma; the same remark immediately records the intended analytic
content: in complex geometry the first stage suffices, and every
non-Archimedean example then known also occurs at the first stage.  CLD call
first-stage functions locally approachable in their Definition~8.1.13.
Thus the coherent open question is early collapse, not eventual
stationarity.  It is distinct from
\Cref{prob:frontier-classical-regularization}: the latter starts with an
arbitrary classically psh function and asks whether it enters the CLD class,
whereas this problem asks how many closure stages are needed after it does.

No example requiring a second stage is known at the cutoff.  A solution must
either flatten an iterated decreasing-limit diagram into a single local
decreasing approximation, with its index and topology controlled, or exhibit
a genuine rank obstruction.  A diagonal argument for sequences is not
enough on a non-first-countable space, and a cofinal sequence may not exist.
The problem remains \emph{open}.

\subsection{Local psh metrics versus global psh metrics}
\index{psh metric!local versus global}

\begin{problem}\label{prob:frontier-cld-global-metric}
Let \(L\) be a line bundle on a projective non-Archimedean variety.  Is every
continuous metric whose local weights are CLD-psh globally psh in the
Fubini--Study or nef-model sense?  More generally, does the corresponding
global approximation statement hold on a proper strict analytic space that
admits a formal model with projective special fibre?
\end{problem}

These are the two questions in \cite[Remark~17.8.10]{CLD25}.  Global
Fubini--Study approximation, semipositive model approximation, and local CLD
positivity agree on the common model loci developed earlier in the book, but
that agreement does not globalize an arbitrary continuous local psh metric.
The missing input is a global approximation theorem respecting transition
functions and the numerical line-bundle class.  Existing global
Fubini--Study and nef-envelope results provide important sufficient cases,
but do not prove the displayed local-to-global implication.  Both questions
remain \emph{open}.

\subsection{Three further Gubler--Jell--Rabinoff questions}
\index{harmonic PL function!semipositive analogue}
\index{abelian variety!canonical metric}

\begin{problem}\label{prob:frontier-gjr-three}
Determine the following.
\begin{enumerate}
\item Does the local finite integral decomposition of harmonic PL functions
      admit a semipositive analogue?
\item For an abelian variety (A), is the canonical injection
      \(\Lambda^{p,q}M_{\mathbf R}\to H^{p,q}(A)\) an isomorphism?
\item Is the canonical metric on an abelian variety an ordinary
      \(\delta\)-metric, rather than only a metric in the larger weakly smooth
      formalism?
\end{enumerate}
\end{problem}

The questions occur, respectively, immediately before
\cite[Proposition~7.5]{GJR26}, in \cite[\S1.15]{GJR26}, and in
\cite[Remark~16.4]{GJR26}.  The last warning is substantive: a previously
used argument relied on smoothness of a canonical tropicalization that fails
in general.  None of the three questions follows from harmonic
tropicalization alone.  They remain \emph{open}; the first is the one most
directly tied to the equivalence problem for local psh definitions.

\section{Valuative singularities and the trivially valued frontier}

\subsection{Do generalized Lelong numbers separate singularities?}

\begin{problem}\label{prob:frontier-lelong-separation}
Let \(u,v\) be psh germs at a smooth complex point.  If
\[
                     \nu_\psi(u)=\nu_\psi(v)
\]
for every tame psh weight \(\psi\), must the valuative transforms
\(\widehat u\) and \(\widehat v\) coincide?
\end{problem}

Boucksom--Favre--Jonsson prove that equality of all infinitely-near Lelong
numbers, equality of strong multiplier ideals, and equality of relative types
are equivalent, and that these conditions imply equality of generalized
Lelong numbers \cite{BFJ08}.  Strong openness later identifies strong and
ordinary multiplier ideals \cite{GZ15}, but it does not reverse the final
implication.  The Zhou-weight results of Bao--Guan--Mi--Yuan
\href{https://doi.org/10.1016/j.aim.2025.110364}{\emph{Adv. Math.} 477
(2025), 110364} strengthen the relative-type and multiplier-ideal testing
language, but do not state this converse for generalized Lelong numbers.

The natural attempted route uses the representation
\eqref{eq:local-generalized-lelong}: do the mixed local Monge--Ampère
measures of tame fundamental solutions separate formal psh functions?
Dimension two and monomial singularities reduce to potential or convex
duality and are the first test cases.  In higher dimension the required
nondegeneracy of the intersection pairing is not presently available in the
needed singular domain.  The converse is \emph{open}.

\subsection{Global Radon--Riesz convergence}

\begin{problem}\label{prob:frontier-bj-radon-riesz}
Does weak convergence together with convergence of the Boucksom--Jonsson
exhaustion \(T_\omega\) imply convergence in the strong
\(d_{\infty,\omega}\)-topology on the full space \(X^{\lin}\)?
\end{problem}

This is the global content of \cite[Question~11.16]{BJ22}.  On a fixed snc
dual complex for a smooth projective variety in characteristic zero, strong
compactness follows from the equi-Lipschitz estimate in
\cite[Theorem~A.4]{BJ22}; compact-to-Hausdorff comparison then identifies the
two topologies, as proved in \Cref{thm:verified-fixed-dual-complex}.  This
argument gives no control uniform in an escaping sequence of models.

The attempted proof seeks a model-independent Izumi estimate and quantitative
control of the retractions.  The counterexample search concentrates on thin
exceptional divisors whose weak coordinates stabilize while their strong
cost escapes.  The fixed-complex case is \emph{verified}; the global
Radon--Riesz statement remains \emph{open}.  Extending the fixed-complex input
to arbitrary characteristic is a separate open compactness problem.

\subsection{A functorial local theory on singular germs}

\begin{problem}\label{prob:frontier-singular-local}
Extend the normalized local theory of
\Cref{chap:local-valuative-psh} from smooth complex germs to excellent
singular and nonnormal germs.  Determine the normalizing boundary datum and
construct, when possible, psh functions, envelopes, bounded products, finite
energy, and finite-map functoriality.
\end{problem}

Several choices that coincide in the smooth one-point setting split here.
The maximal ideal, a boundary divisor, and a chosen ideal of definition can
cut different links in the valuation cone.  Normalization separates
branches; descent must remember the conductor.  In positive characteristic,
alterations replace resolutions but do not automatically carry the local
intersection product required by \Cref{thm:local-formal-ma}.

The Rethlas programme tests normal crossings, cusps, nodes, quotient
singularities, and purely inseparable universal homeomorphisms.  Any claim of
choice-independence must exhibit its comparison maps and scaling law.
Relative b-divisors can be defined on suitable model systems, but the full
package above is presently \emph{open}.

\index{local volume!resolved questions}
Two older local-volume questions should not be carried forward as open.
For the Boucksom--de Fernex--Favre local volume, Yuchen Zhang proves a
positive dimension-dependent lower bound for every nonzero volume of an
isolated Gorenstein singularity and constructs a non-
\(\mathbf Q\)-Gorenstein isolated singularity of volume zero that admits no
boundary making the pair log canonical
\cite{Zhang14LocalVolume}.  Thus the first conjectural
direction is confirmed and the second implication is disproved.  These
results delimit the singular theory above; they are resolved history, not
additional problems.

\subsection{Finite-energy equations on polarized affine cones}

\begin{problem}\label{prob:frontier-affine-cone-ma}
Let \((X,T;\xi)\) be a normal polarized affine cone over an algebraically
closed trivially valued field of characteristic zero.  Develop a comparison
calculus for \(\PSH(\xi)\) and \(\PSH(\xi')\) when \(\xi\) and \(\xi'\) are
nonproportional Reeb fields, strong enough to prove the required uniform
differentiability of the energy.  For a fixed \(\xi\), characterize the image
of
\[
 \MA_\xi:\mathcal E^1(\xi)/\mathbf R
 \longrightarrow \mathcal M^1(\xi),
\]
and determine exactly which \(\mu\in\mathcal M^1(\xi)\) admit a solution of
\(\MA_\xi(\varphi)=\mu\), unique modulo constants.
\end{problem}

The fixed-polarization theory is already substantial:
\Cref{thm:local-cone-finite-energy-ma} constructs the mixed operator,
proves decreasing-net continuity, and makes the single-potential operator
bi-Lipschitz onto its image.  The proportional-rescaling identities of
\cite[Proposition~3.3]{Wu24AffineCones} do not compare nonproportional Reeb
rays.  Rational approximation likewise does not yield constants uniform in
two varying polarizations.  The introduction of \cite{Wu24AffineCones}
identifies precisely this missing uniform differentiability input and does
not claim surjectivity onto \(\mathcal M^1(\xi)\).

The natural strategy is to prove Reeb-uniform Hodge-index and
\(I\)--\(J\) estimates on a compact slice of the Reeb cone, use them to
differentiate the energy along rational approximants, and then run the
finite-energy variational method.  At present the first step is unavailable:
the known addition formula stays on one Reeb ray, and an irrational output
measure may already be carried by a genuinely quasimonomial point.  No
argument in this edition closes either the cross-polarization estimate or the
image characterization.  Both residuals are therefore \emph{open} and are
recorded as one coupled, unresolved Rethlas packet in the editorial ledger;
no verification verdict is claimed for either residual.

\subsection{The formerly listed trace refinement is resolved}
\label{sec:frontier-trace-p-resolved}

An earlier draft asked whether analytic trace on its natural domain admits
a map
\[
 D_Y/{\sim_P}\longrightarrow
 \QPSH(\widetilde Y)/{\sim_P}.
\]
This is not open.  The trace theorem proves the stronger factorization
\[
 D_Y/{\sim_P}
 \longrightarrow
 D_Y/{\sim_{\mathcal I}}
 \longrightarrow
 \QPSH(\widetilde Y)/{\sim_P}:
 \qquad
 [\varphi]_{\mathcal I}\longmapsto[\Tr_Y(\varphi)]_P.
\]
Every trace representative is \(\mathcal I\)-good; moreover
\(\varphi\preceq_{\mathcal I}\psi\) implies
\(\Tr_Y(\varphi)\preceq_P\Tr_Y(\psi)\).  Thus
\(\mathcal I\)-equivalent inputs have \(P\)-equivalent traces.  These are
the comparison and uniqueness statements of \cite{DX24}, reproduced with
full proofs in the supplied trace chapter of \cite{XiaBook}; see
\Cref{thm:operations-trace-well-defined,thm:operations-trace-p-descent,thm:verified-trace-factorization}.

This result does not assert that the input relations
\(\sim_{\mathcal I}\) and \(\sim_P\) coincide.  It says instead that the
restriction construction improves the target uniqueness.  The old problem
environment has therefore been removed from the active inventory rather
than relabelled as a new result of this book.

\section{Pseudoeffective classes and zero volume}

\subsection{Extending the MP realization to every pseudoeffective class}

\begin{problem}\label{prob:frontier-mp-pseff}
Let \(X\) be a compact Kähler manifold and let
\(\alpha\in H^{1,1}_{\mathrm{BC}}(X,\mathbf R)\) be pseudoeffective.
Construct the strongest intrinsic realization on the Mesquita--Piccione
tropical spectrum of the DXZ inverse-system space
\(\PSH^{\NA}(X,\alpha)\), and identify its order, homogeneity,
approximation, energy, and Monge--Ampère structures whenever these exist.
\end{problem}

The verified core is now \Cref{chap:pseff-mp-dxz}.  Every compatible DXZ
family defines a perturbation-independent raw ideal-valuative function on
\(X^\beth\).  Its divisorial restriction is the sourced DXZ
analytification, which is injective and order-reflecting; no upper
semicontinuity or MP positivity is inferred at an arbitrary nondivisorial
point on a general compact K\"ahler manifold.  Translation holds on all
semivaluations, while the other pointwise operation formulas are retained
on the rational divisorial locus with the exact endpoint hypotheses of
\cite{Xia23Operations}.  The K\"ahler finite-energy correspondence with MP
is proved.  If \(X\) is smooth projective and the class is pseudoeffective
in the real N\'eron--Severi space, the full DXZ--Boucksom--Jonsson
correspondence is known; for a \(\mathbf Q\)-line-bundle class, the raw
ideal-valuative function is exactly the Boucksom--Jonsson psh metric on
every semivaluation by \cite[formula~(3.8) and Theorem~3.12]{Xia23Operations}.
Even for an ample \(\mathbf Q\)-class, however, unrestricted identification
with the MP cone remains open.

The problem remains open in three precise senses: characterize the image of
the transform intrinsically; decide whether the raw transform is upper
semicontinuous or MP psh at arbitrary nondivisorial points on a general
compact K\"ahler manifold; and construct a \(P\)-sensitive positivity,
measure, and energy theory for big non-nef classes.  One cannot obtain the
latter by replacing \(P\)-equivalence with \(\mathcal I\)-equivalence.

There is also a sharper published big-class question hidden inside the first
item.  Darvas--Xia--Zhang conjecture, immediately before their
Definition~4.4, that filtration test curves are maximal in the big case
\cite{DXZ23}.  Proving this in the positive-mass range would identify an
intrinsic part of the image of the realization and remove one envelope
operation from the comparison with MP.  The zero-mass range also contains
the separate open problem of recovering an \(\mathcal I\)-model potential
from data with vanishing top mass, so it must not be treated as a formal
endpoint of the positive-mass argument.  Both ranges remain \emph{open};
the candidate-only audit in the editorial ledger produced no Rethlas
receipt and hence no promoted endpoint theorem.

Two unchanged extensions are excluded by
\Cref{prop:pseff-relative-nef-obstruction,prop:pseff-zero-volume-boundary}.
A relatively \emph{nef} total-space class restricts to a nef class on every
noncentral fibre, so vertical corrections and central modifications that
remain relatively nef cannot represent a big non-nef generic class.  If
\(\operatorname{vol}(\alpha)=0\), every top nonpluripolar mass vanishes and
the existing probability and normalized energy formulas are undefined.

\subsection{Zero-volume energy, measures, and convex bodies}

\begin{problem}\label{prob:frontier-zero-volume}
After choosing the minimum necessary reference scale, is there a
path-independent unnormalized energy and Monge--Ampère theory for a
pseudoeffective zero-volume class?  Is there a \(P\)-sensitive
transcendental Okounkov body encoding its lower-dimensional numerical data?
\end{problem}

The unchanged probability signature is \emph{disproved}: total mass zero
cannot be normalized to mass one.  Kähler perturbations give a family of
probabilities, not automatically a canonical limit.  Moreover
\Cref{thm:verified-mass-no-go} proves that a body depending only on
\(\mathcal I\)-equivalence cannot recover actual nonpluripolar mass.

The attempted positive construction rescales by the first nonzero coefficient
in \(\operatorname{vol}(\alpha+\varepsilon\omega)\), then tests independence
of the perturbation path and comparison with numerical dimension.  Nef
zero-volume classes and toric examples provide the first nontrivial cases.
Any surviving convex body must use \(P\)-type or richer data.  No general
path-independent construction is presently verified.

Two source questions give sharper benchmarks.  Xia conjectures that the
dimension of the partial Okounkov body equals the numerical dimension even
at zero mass \cite[Conjecture~5.29]{Xia21}, and the same source asks whether
the Chebyshev transform is independent of its auxiliary positive form
\cite[p.~34]{Xia21}.  These remain genuine published questions.  The
\(P\)-sensitive body requested above is a stronger construction formulated
in this book and should not be attributed to that source.

The third conjecture of that source is the non-toric mixed-volume formula
\cite[Conjecture~8.8]{Xia21}.  On 24 August 2026, the author supplied the TeX
proof underlying the unpublished manuscript \cite{XiaMVOB}.  After repairs
of its positive-Cao-volume reduction, convex sandwich, pointwise potential
sandwich, rational analytic-resolution step, degenerate finite-level case,
multigraded saturation step, and effective-shift body inclusion, the exact
repaired proof received a production \texttt{gpt-5.6-sol}/\texttt{xhigh}
verdict of \emph{correct}, with no critical error and no gap.  Thus
\Cref{thm:ok-mvob-metrics} proves Conjecture~8.8 for singular psh metrics on
big line bundles over a smooth projective complex variety when every Cao
(equivalently, ideal-model) volume is positive;
\Cref{thm:ok-mvob-line-bundles} gives the movable-product specialization for
big line bundles with minimal singularities.  This is an author-supplied,
repaired, and independently verified proof in an unpublished work in
progress, not a claim that the result has appeared in the published
literature.  It asserts a supremum over surjective rank-\(n\) valuations,
not attainment by one valuation.  Zero volume, transcendental classes,
singular varieties, and attainment remain outside the theorem; in
particular, the genuinely different problem
\Cref{prob:frontier-zero-volume} remains open.
This source-status item is not an additional problem environment in the
literal count at the beginning of the chapter.

\section{Existence, degenerations, and numerical dimension}

\subsection{The Monge--Ampère equation beyond the original BFJ domain}

\begin{problem}\label{prob:frontier-general-ma}
Extend the finite-energy existence and uniqueness theory for the NA
Monge--Ampère equation beyond the currently proved combinations of smoothness,
projectivity, residue characteristic, valuation, positivity class, support,
and algebraizability hypotheses.
\end{problem}

This is not one yes-or-no assertion.  Positive characteristic, singular or
unibranch varieties, nondiscrete valuations, big classes, arbitrary support,
and nonalgebraizable analytic spaces require different inputs.  The known
variational proofs use compactness, envelope regularity, differentiability of
energy, and a comparison principle; deleting a geometric hypothesis is valid
only after replacing every place where it enters.  The relevant positive
results and their distinct domains are surveyed in
\Cref{chap:bfj-finite-energy-measures,chap:bfj-continuity}.

One exact source-stated axis is especially clean.  Pille--Schneider observes
that the finite-energy space can be defined over an arbitrary
non-Archimedean field, but that the extension of the Monge--Amp\`ere operator
is not known there in general; the discretely valued characteristic-zero and
trivially valued cases are the established endpoints
\cite[Remark~4.4]{PS23Hybrid}.  This should be solved before bundling all
remaining axes into a single theorem.

\index{Monge--Amp\`ere equation!skeletal factorization}
There is also a resolved negative regularity chronology.  The first
skeletal-factorization question of Gubler--Jell--K\"unnemann--Martin already
fails on their toric surface \cite[Appendix~A]{GJKM19}.
Boucksom--Jonsson then construct a finite
atomic Monge--Amp\`ere measure supported on a dual complex whose solution is
invariant under no retraction attached to any SNC model; they also construct
a divisorial Green function that is not \(\mathbf Q\)-PL
\cite[Theorems~A and~B]{BJGreen24}.  Hence support on a dual complex does not
imply factorization through some refined skeleton.  This is a theorem-level
obstruction to a proof strategy, not an open regularity expectation.

The Rethlas programme treats each axis separately and first searches the
post-2020 literature for an existing theorem.  Alterations and universal
homeomorphisms can address some algebraic descent questions, but they do not
by themselves supply multiplier ideals or differentiability in positive
characteristic.  The unrestricted formulation is \emph{open}; several
strictly delimited axes are already known.

\subsection{Continuity of Monge--Amp\`ere measures over Banach bases}

\begin{problem}\label{prob:frontier-banach-family-ma}
Let \(A\) be an integral Banach ring, let \(X\) be a flat projective
\(A\)-scheme of relative dimension \(n\), let \(L_1,\ldots,L_n\) be
semiample line bundles, and let \(\phi_i\) be continuous psh metrics on
\(L_i\).  For every \(\psi\in C^0(X^{\mathrm{an}})\), is
\[
 x\longmapsto
 \int_{X_{\mathcal H(x)}^{\mathrm{an}}}\psi\,
 \operatorname{MA}\bigl((\phi_1)_x,\ldots,(\phi_n)_x\bigr)
\]
continuous on \(\mathcal M(A)\)?
\end{problem}

This is \cite[Question~4.9]{PS23Hybrid}.  Flatness is essential even to keep
the total
mixed intersection mass constant.  The answer is affirmative for the hybrid
base and continuous hybrid psh metrics by
\cite[Theorem~4.11]{PS23Hybrid}, which reformulates Favre's theorem.  The
source states that the answer is unclear for a
general integral Banach ring and expects the arithmetic case of the ring of
integers of a number field.  Its Chern--Levine--Nirenberg estimate
\cite[Lemma~4.10]{PS23Hybrid} and density argument
reduce the question to FS metrics and model test functions.  The Rethlas task
is to prove that reduced assertion by an intersection-theoretic
specialization formula, or to produce a non-hybrid Banach base where
specialization of vertical intersection cycles jumps.  The general problem
is \emph{open}.

\subsection{Continuity of the hybrid extension}

\begin{problem}\label{prob:frontier-favre-hybrid-extension}
Let \((X,L)\to\mathbb D^*\) be a polarized projective degeneration.  Let
\(\phi\) be a continuous semipositive metric with logarithmic growth and
assume that its associated non-Archimedean metric \(\phi^{\NA}\) is
continuous.  Is the induced metric \(\phi^{\mathrm{hyb}}\) continuous on the hybrid
space, and hence a continuous hybrid psh metric?
\end{problem}

This is Favre's Question~1 as reproduced in
\cite[Question~3.18]{PS23Hybrid}.  The psh property of the glued metric is
\cite[Theorem~3.13]{PS23Hybrid}; continuity is the missing assertion.  It is
proved for degenerations of canonically polarized manifolds admitting
semistable reduction by \cite[Theorem~5.1]{PS23Hybrid}, and continuous hybrid
metrics have weakly continuous Monge--Amp\`ere measures by its Theorem~4.11.
The proposed proof requires a Bergman
regularization error uniform in the degeneration parameter, equivalently the
uniform oscillation estimate isolated immediately after the question in the
primary source.  Continuity of the two endpoint metrics alone does not give
that estimate.  The general problem remains \emph{open}.

\subsection{Hybrid entropy slopes and the remaining compact-Kähler YTD frontier}

\begin{problem}\label{prob:frontier-hybrid-ytd}
Prove the entropy-slope, approximation, compactness, and automorphism
statements needed to pass from model non-Archimedean stability to existence
of canonical metrics along arbitrary finite-energy weak rays on a compact
Kähler manifold.
\end{problem}

The static projective entropy and stability gap must not be listed as wholly
open at this cutoff.
Trusiani's version-two preprint of 17 August 2026,
\cite[Theorem~B and Corollary~A]{Trusiani26}, retains the special Fujita
approximations for big line bundles and the claimed solution of the
Boucksom--Jonsson entropy-regularization conjecture; its Theorem~A retains
the automorphism-relative uniform YTD theorem for smooth polarized
projective varieties.  Corollary~A also records the resulting equality of
the Archimedean entropy slope along a maximal finite-energy ray with the
non-Archimedean entropy, in that projective framework.  The new Appendix by
S\'ebastien Boucksom supplies a non-Archimedean proof over
\(\Bbbk((t))\), and Theorem~C adds the weighted-relative correspondence,
including Calabi extremal metrics when the weight is log-concave.  These are
recent preprint claims, rather than results independently reproved in this
book.  Boucksom--Jonsson's 18 September 2025
preprint \cite{BJ25WeightedYTD} gives a
non-Archimedean formulation for cscK and weighted extremal metrics on
polarized projective manifolds.  Darvas--Zhang's 3 November 2025 version-two
preprint \cite[Theorems~1.6 and~5.5]{DarvasZhang25YTD} gives a uniform
YTD correspondence for a unique cscK metric on an arbitrary compact
Kähler manifold using \(K^\beta\)-stability and, in the polarized
projective case, a reductivity-plus-log-discrepancy-model formulation.
Mesquita--Piccione--Witt Nyström prove that uniform model stability implies
existence in a Kähler class and establish the transcendental finite-energy
Calabi--Yau homeomorphism \cite[Theorems~A and~B]{MPWN25}.  Finally,
Trusiani's Corollary~B identifies the standard automorphism-relative
projective stability formulations in the smooth polarized setting
\cite{Trusiani26}.  Thus the projective entropy slope, the weighted-relative
correspondence in the stated smooth polarized scope, and comparison of the
listed projective signatures are no longer open problems.

\Cref{chap:hybrid-spaces,chap:ytd-bridge} separates hybrid convergence from
entropy and coercivity.  The BHJ erratum is enforced in
\Cref{chap:rays-filtrations}.  The remaining targets are narrower: direct
algebraic proofs of the known projective equivalences; \(K^\beta\)-stability
in the presence of nontrivial automorphisms; equivariant transcendental and
singular extensions; regularization of rays attached to log-discrepancy
models; and the genuinely hybrid compactness and measure-continuity
statements needed to pass through a degenerating family.  Boucksom--Jonsson
also isolate bounded support size for divisorial test measures, algebraic
equivariance proofs, and extension to polarized klt pairs
\cite[Questions and further remarks]{BJ25WeightedYTD}.  Those residual
problems remain \emph{open}; the smooth polarized projective entropy and
weighted-relative theorems are dated Trusiani version-two claims, not
missing inputs.

\subsection{Irrational real toric numerical dimension}

\begin{problem}\label{prob:frontier-real-toric}
For a torus-invariant positive current \(T\) in a big irrational real
class on a smooth projective toric variety, does
\[
                         \nu(T)=\dim\Delta(T)
\]
hold for the numerical dimension defined by real positive products?
\end{problem}

\Cref{thm:verified-toric-dimension} proves the identity for integral and
rational toric classes using Guenancia's multiplier-ideal criterion
\cite[Theorem~1.13]{Guenancia12}, the character decomposition
\cite[\S4.3]{CLS11}, and Cao's projective numerical criterion
\cite[Proposition~4.4]{Cao14}.  A real class has no single section algebra to
which that argument applies.

The attempted extension compares positive products under rational
class-varying approximation and seeks a purely convex characterization of
the transcendental numerical dimension.  Uniform control of singularity type
is essential; convergence of Newton bodies alone does not identify multiplier
ideals or positive products.  The irrational real statement is
\emph{open}.

\subsection{Realizing the cross-fan counterexample analytically}

\begin{problem}\label{prob:frontier-cross-fan}
Can the weighted cross-fan of
\Cref{prop:verified-cross-fan} be realized by a harmonic tropicalization of a
good strict analytic space and by a positive closed analytic or delta-current
with the displayed weights?
\end{problem}

The abstract calibrated fan counterexample is verified, but analytic
realization requires more than balancing.  A candidate must construct the
space, tropical chart, lattice calibration, multiplicities, and current, and
must show that the ambient convex approximants pull back to the asserted psh
functions.  Nodal curves, products of annuli, semistable surface skeleta, and
tropical modifications are the initial tests.  Conversely, analytic
multiplicity formulas might force the current weights to be proportional to
the fundamental calibration, which would be a negative realization theorem.
No such construction or obstruction is presently verified.

\section{How solutions enter the book}

Each problem above is assigned a separate, source-scoped Rethlas input so
that a solution of one question cannot hide a gap in another.  A queued input
is not described as an attempted solution until a final verifier receipt
exists.  A successful proof is first checked
against its primary sources and boundary cases, then verified in a fresh
session, and finally reread after integration.  A counterexample is held to
the same standard: abstract tropical data, analytic realizations, and global
algebraic examples are not interchangeable.

The current results of this process are already visible:

\begin{enumerate}
\item the literal equality of smooth, PL, current, classical, and global psh
      classes is disproved by elementary regularity and polarization
      obstructions;
\item the three PL notions in
      \eqref{eq:verified-pl-psh-equivalence} are verified to coincide on the
      relative interior;
\item restriction of a positive aggregate tropical divisor to a balanced
      curve is disproved in general, while the corresponding analytic CLD
      converse remains open conditional on
      \eqref{eq:frontier-restriction-r};
\item probability-valued local restriction is disproved as a sheaf
      signature;
\item unchanged MP models are disproved for a non-Kähler generic class, and
      unchanged normalized energy is disproved at volume zero;
\item the \(\mathcal I\)-invariant mass-body proposal is disproved by
      \Cref{thm:verified-mass-no-go};
\item support of the right-hand side on a dual complex is disproved as a
      criterion for invariance under some SNC retraction
      \cite{BJGreen24};
\item the two older BdFF local-volume questions are respectively confirmed
      and disproved by \cite{Zhang14LocalVolume}; and
\item the fixed-dual-complex, split-trace, rational toric dimension, and
      cross-fan statements are verified in their displayed domains.
\end{enumerate}

The unresolved cores remain problems, not promises.  Their attempted routes
and failed implications are retained so that future revisions can update one
status at a time without changing the surrounding theory.

\section*{Notes and further reading}

The CLD and Gubler--Rabinoff boundaries use
\cite{CLD25,GJR26,GR25PL,GR25Classical,Wanner19}; the polyhedral frontier is
\href{https://arxiv.org/abs/2603.08664}{arXiv:2603.08664v1}.  The local valuative
questions originate in \cite{BFJ08}; the global strong-topology question is
\cite[Question~11.16]{BJ22}.  The pseudoeffective inverse system is due to
\cite{DXZ23}, with operations developed in \cite{Xia23Operations}; the
Kähler tropical-spectrum theory is \cite{MP24}.  The two exact hybrid questions
are recorded in \cite{PS23Hybrid}.  Stability status at the cutoff uses
\cite{Trusiani26,BJ25WeightedYTD,DarvasZhang25YTD}, with their
different signatures kept separate.  Further routing appears in
\Cref{chap:hybrid-spaces,chap:ytd-bridge}.  Because this literature is rapidly
changing, every status in this chapter carries the cutoff date printed at its
beginning.
